% =====================================================================
%  Appendix: Appendix (A)
%  Gattringer & Lang, "Quantum Chromodynamics on the Lattice" (LNP 788)
%  Reconstructed from extract/ch13.txt; equations auto-numbered.
%  main.tex issues \appendix before % =====================================================================
%  Appendix: Appendix (A)
%  Gattringer & Lang, "Quantum Chromodynamics on the Lattice" (LNP 788)
%  Reconstructed from extract/ch13.txt; equations auto-numbered.
%  main.tex issues \appendix before % =====================================================================
%  Appendix: Appendix (A)
%  Gattringer & Lang, "Quantum Chromodynamics on the Lattice" (LNP 788)
%  Reconstructed from extract/ch13.txt; equations auto-numbered.
%  main.tex issues \appendix before % =====================================================================
%  Appendix: Appendix (A)
%  Gattringer & Lang, "Quantum Chromodynamics on the Lattice" (LNP 788)
%  Reconstructed from extract/ch13.txt; equations auto-numbered.
%  main.tex issues \appendix before \input{appendix}, so the chapter is
%  automatically lettered "A" and the sections become A.1, A.2, ...
% =====================================================================

\chapter{Appendix}

\section{The Lie groups SU($N$)}

In this appendix we collect basic definitions and conventions for the Lie
groups SU($N$) -- the special unitary groups -- and the corresponding Lie
algebras $su(N)$. For a more detailed presentation we refer the reader to
[1--3].

\subsection{Basic properties}

The defining representation of SU($N$) is given by complex $N\times N$
matrices which are unitary and have determinant 1. This set of matrices is
closed under matrix multiplication: Let $\Omega_1$ and $\Omega_2$ be elements
of SU($N$), i.e., they obey $\Omega_i^\dagger = \Omega_i^{-1}$ and
$\det[\Omega_i] = 1$. Using standard linear algebra manipulations we obtain
\begin{equation}
  \begin{gathered}
    (\Omega_1\Omega_2)^\dagger = \Omega_2^\dagger\Omega_1^\dagger
    = \Omega_2^{-1}\Omega_1^{-1} = (\Omega_1\Omega_2)^{-1},\\[2pt]
    \det[\Omega_1\Omega_2] = \det[\Omega_1]\det[\Omega_2] = 1
  \end{gathered}
\end{equation}
and have thus established that also the product of two SU($N$) matrices is an
SU($N$) matrix. The unit matrix is also in SU($N$) and for each matrix in
SU($N$) there exists an inverse (the hermitian conjugate matrix). Thus, the
set SU($N$) forms a group. Since the group operation -- the matrix
multiplication -- is not commutative the groups SU($N$) are non-abelian
groups.

\subsection{Lie algebra}

Let us now count how many real parameters are needed to describe the matrices
in SU($N$). A complex $N\times N$ matrix has $2N^2$ real parameters. The
requirement of unitarity introduces $N^2$ independent conditions which the
parameters have to obey. One more parameter is used for obeying the determinant
condition such that one needs a total of $N^2-1$ real parameters for describing
SU($N$) matrices.

A convenient way of representing SU($N$) matrices is to write them as
exponentials of basis matrices $T_j$, the so-called generators. In particular,
we write an element $\Omega$ of SU($N$) as
\begin{equation}
  \Omega = \exp\!\left(i\sum_{j=1}^{N^2-1}\omega^{(j)} T_j\right),
\end{equation}
where $\omega^{(j)}$, $j = 1, 2, \ldots, N^2-1$, are the real numbers needed to
parameterize $\Omega$. We remark that the parameters $\omega^{(j)}$ can be
changed continuously, making SU($N$) so-called Lie groups that are groups whose
elements depend continuously on their parameters. In order to cover all of the
group space, the parameters have to be varied only over finite intervals,
making the groups SU($N$) so-called compact Lie groups.

The generators $T_j$, $j = 1, 2, \ldots, N^2-1$, are chosen as traceless,
complex, and hermitian $N\times N$ matrices obeying the normalization condition
\begin{equation}
  \tr[T_j T_k] = \half\,\delta_{jk}.
\end{equation}
In addition, they are related among each other by an algebra of commutation
relations
\begin{equation}
  [T_j, T_k] = i\, f_{jkl}\, T_l .
\end{equation}
The completely anti-symmetric coefficients $f_{jkl}$ are the so-called
structure constants. Below we will give an explicit representation of the
generators for the groups SU(2) and SU(3).

Let us verify that the representation of $\Omega$ indeed describes elements of
SU($N$). Using the facts that the generators are hermitian and that the
$\omega^{(j)}$ are real, one finds that hermitian conjugation of the right-hand
side of the exponent simply produces an extra minus sign in the exponent (from
the complex conjugation of $i$). Thus, it is obvious that $\Omega^\dagger =
\Omega^{-1}$. To show that the determinant equals 1, we use the equation
\begin{equation}
  \det\Omega = \exp(\tr[\ln\Omega])
  = \exp\!\left(i\sum_{j=1}^{N^2-1}\omega^{(j)}\tr T_j\right) = e^0 = 1,
\end{equation}
where in first step we have used a formula for the determinant (see (A.54)
below) and in the third step we have used the fact that the $T_j$ are
traceless.

Not only the group elements but also the exponents of our representation have
an interesting structure. The linear combinations
\begin{equation}
  \sum_{j=1}^{N^2-1}\omega^{(j)} T_j
\end{equation}
of the $T_j$ form the so-called Lie algebra $su(N)$. Their commutation
properties are governed by the commutation relations. Elements of $su(N)$ are
also complex $N\times N$ matrices but have properties different from the
elements of the group. One important difference is the fact that the unit
matrix is contained in the group (all $\omega^{(j)} = 0$), while it is not an
element of the algebra (all $T_j$ are traceless).

\subsection{Generators for SU(2) and SU(3)}

The standard representation of the generators for SU(2) is given by
\begin{equation}
  T_j = \half\,\sigma_j,
\end{equation}
with the Pauli matrices
\begin{equation}
  \sigma_1 = \begin{pmatrix} 0 & 1 \\ 1 & 0 \end{pmatrix},\qquad
  \sigma_2 = \begin{pmatrix} 0 & -i \\ i & 0 \end{pmatrix},\qquad
  \sigma_3 = \begin{pmatrix} 1 & 0 \\ 0 & -1 \end{pmatrix}.
\end{equation}
In this case the structure constants are particularly simple, given by the
completely anti-symmetric tensor, i.e., $f_{jkl} = \varepsilon_{jkl}$.

For SU(3) the generators are given by
\begin{equation}
  T_j = \half\,\lambda_j .
\end{equation}
The Gell--Mann matrices $\lambda_j$ are $3\times 3$ generalizations of the
Pauli matrices:
\begin{equation}
  \begin{gathered}
    \lambda_1 = \begin{pmatrix} 0 & 1 & 0 \\ 1 & 0 & 0 \\ 0 & 0 & 0 \end{pmatrix},\qquad
    \lambda_2 = \begin{pmatrix} 0 & -i & 0 \\ i & 0 & 0 \\ 0 & 0 & 0 \end{pmatrix},\qquad
    \lambda_3 = \begin{pmatrix} 1 & 0 & 0 \\ 0 & -1 & 0 \\ 0 & 0 & 0 \end{pmatrix},\\[4pt]
    \lambda_4 = \begin{pmatrix} 0 & 0 & 1 \\ 0 & 0 & 0 \\ 1 & 0 & 0 \end{pmatrix},\qquad
    \lambda_5 = \begin{pmatrix} 0 & 0 & -i \\ 0 & 0 & 0 \\ i & 0 & 0 \end{pmatrix},\qquad
    \lambda_6 = \begin{pmatrix} 0 & 0 & 0 \\ 0 & 0 & 1 \\ 0 & 1 & 0 \end{pmatrix},\\[4pt]
    \lambda_7 = \begin{pmatrix} 0 & 0 & 0 \\ 0 & 0 & -i \\ 0 & i & 0 \end{pmatrix},\qquad
    \lambda_8 = \frac{1}{\sqrt{3}}
    \begin{pmatrix} 1 & 0 & 0 \\ 0 & 1 & 0 \\ 0 & 0 & -2 \end{pmatrix}.
  \end{gathered}
\end{equation}

\subsection{Derivatives of group elements}

Let us now show an important property of derivatives of group elements. If
$\Omega(\omega)$ is an element of SU($N$) then
\begin{equation}
  M_k(\omega) = i\,\frac{\partial\Omega(\omega)}{\partial\omega^{(k)}}\,
  \Omega(\omega)^\dagger \in su(N),
\end{equation}
i.e., the derivative times the hermitian conjugate is in the Lie algebra. In
order to prove this statement we have to show the defining properties of Lie
algebra elements, i.e., we must show that $M_k(\omega)$ is hermitian and
traceless.

Showing the hermiticity of $M_k(\omega)$ is straightforward. By differentiating
$\Omega(\omega)\Omega(\omega)^\dagger = 1$ with respect to $\omega^{(k)}$ one
finds
\begin{equation}
  \frac{\partial\Omega(\omega)}{\partial\omega^{(k)}}\,\Omega(\omega)^\dagger
  + \Omega(\omega)\,\frac{\partial\Omega(\omega)^\dagger}{\partial\omega^{(k)}} = 0.
\end{equation}
Thus
\begin{align}
  M_k(\omega)^\dagger &=
  \left[ i\,\frac{\partial\Omega(\omega)}{\partial\omega^{(k)}}\,
  \Omega(\omega)^\dagger \right]^\dagger
  = -i\,\Omega(\omega)\,\frac{\partial\Omega(\omega)^\dagger}{\partial\omega^{(k)}} \notag\\
  &= i\,\frac{\partial\Omega(\omega)}{\partial\omega^{(k)}}\,
  \Omega(\omega)^\dagger = M_k(\omega),
\end{align}
where we used (A.12) in the third step.

In order to show that $M_k(\omega)$ is traceless we use the fact that the
determinant of an SU($N$) matrix equals 1 and we differentiate
$\det[\Omega(\omega)]$ with respect to $\omega^{(k)}$. We obtain
\begin{align}
  0 &= \frac{\partial\det[\Omega(\omega)]}{\partial\omega^{(k)}}
  = \frac{\partial\det[\Omega(\omega)]}{\partial\Omega(\omega)_{ab}}\,
  \frac{\partial\Omega(\omega)_{ab}}{\partial\omega^{(k)}} \notag\\
  &= \det[\Omega(\omega)]\left(\Omega(\omega)^{-1}\right)_{ba}\,
  \frac{\partial\Omega(\omega)_{ab}}{\partial\omega^{(k)}}
  = \tr\!\left[\frac{\partial\Omega(\omega)}{\partial\omega^{(k)}}\,
  \Omega(\omega)^\dagger\right],
\end{align}
where in the first step we applied the chain rule for derivatives. In the second
step we used a standard formula for the derivative of the determinant
$\det[\Omega]$ with respect to an entry $\Omega_{ab}$ of the matrix $\Omega$.
In the third step we used $\det[\Omega] = 1$ and $\Omega^{-1} = \Omega^\dagger$.
Equation (A.14) establishes that $M_k(\omega)$ is also traceless and thus we
have shown $M_k(\omega)\in su(N)$.

From (A.11) it follows that for the gauge transformation matrices $\Omega(x)$
with coefficients $\omega^{(k)}(x)$, depending on the space--time coordinate
$x$, the combination $i(\partial_\mu\Omega(x))\Omega(x)^\dagger$ is in the Lie
algebra, since
\begin{equation}
  i(\partial_\mu\Omega(x))\,\Omega(x)^\dagger
  = \sum_k i\,\frac{\partial\Omega(\omega(x))}{\partial\omega^{(k)}(x)}\,
  \Omega(\omega(x))^\dagger\,\partial_\mu\omega^{(k)}(x),
\end{equation}
and the right-hand side is a linear combination of $su(N)$ elements with real
coefficients $\partial_\mu\omega^{(k)}(x)$.

\section{Gamma matrices}

The Euclidean gamma matrices $\gamma_\mu$, $\mu = 1, 2, 3, 4$ can be
constructed from the Minkowski gamma matrices $\gamma_\mu^M$, $\mu = 0, 1, 2,
3$. The latter obey
\begin{equation}
  \{\gamma_\mu^M, \gamma_\nu^M\} = 2\,g_{\mu\nu}\,\mathbb{1},
\end{equation}
with the metric tensor given by $g_{\mu\nu} = \mathrm{diag}(1,-1,-1,-1)$ and
$\mathbb{1}$ is the $4\times 4$ unit matrix. Thus when we define the Euclidean
matrices $\gamma_\mu$ by setting
\begin{equation}
  \gamma_1 = -i\gamma_1^M,\quad
  \gamma_2 = -i\gamma_2^M,\quad
  \gamma_3 = -i\gamma_3^M,\quad
  \gamma_4 = \gamma_0^M,
\end{equation}
we obtain the Euclidean anti-commutation relations
\begin{equation}
  \{\gamma_\mu, \gamma_\nu\} = 2\,\delta_{\mu\nu}\,\mathbb{1}.
\end{equation}

In addition to the matrices $\gamma_\mu$, $\mu = 1, 2, 3, 4$ we define the
matrix $\gamma_5$ as the product
\begin{equation}
  \gamma_5 = \gamma_1\gamma_2\gamma_3\gamma_4 .
\end{equation}
The matrix $\gamma_5$ anti-commutes with all other gamma matrices $\gamma_\mu$,
$\mu = 1, 2, 3, 4$ and obeys $\gamma_5^2 = 1$.

An explicit representation of the Euclidean gamma matrices can be obtained from
a representation of the Minkowski gamma matrices (see, e.g., [4]). Here we give
the so-called chiral representation where $\gamma_5$ (the chirality operator)
is diagonal:
\begin{equation}
  \gamma_{1,2,3} = \begin{pmatrix} 0 & -i\sigma_{1,2,3} \\ i\sigma_{1,2,3} & 0 \end{pmatrix},\qquad
  \gamma_4 = \begin{pmatrix} 0 & \mathbb{1}_2 \\ \mathbb{1}_2 & 0 \end{pmatrix},\qquad
  \gamma_5 = \begin{pmatrix} \mathbb{1}_2 & 0 \\ 0 & -\mathbb{1}_2 \end{pmatrix},
\end{equation}
where the $\sigma_j$ are the Pauli matrices and $\mathbb{1}_2$ is the
$2\times 2$ unit matrix. More explicitly the Euclidean gamma matrices read
\begin{equation}
  \begin{gathered}
    \gamma_1 = \begin{pmatrix} 0 & 0 & 0 & -i \\ 0 & 0 & -i & 0 \\ 0 & i & 0 & 0 \\ i & 0 & 0 & 0 \end{pmatrix},\qquad
    \gamma_2 = \begin{pmatrix} 0 & 0 & 0 & -1 \\ 0 & 0 & 1 & 0 \\ 0 & 1 & 0 & 0 \\ -1 & 0 & 0 & 0 \end{pmatrix},\\[4pt]
    \gamma_3 = \begin{pmatrix} 0 & 0 & -i & 0 \\ 0 & 0 & 0 & i \\ i & 0 & 0 & 0 \\ 0 & -i & 0 & 0 \end{pmatrix},\qquad
    \gamma_4 = \begin{pmatrix} 0 & 0 & 1 & 0 \\ 0 & 0 & 0 & 1 \\ 1 & 0 & 0 & 0 \\ 0 & 1 & 0 & 0 \end{pmatrix},\\[4pt]
    \gamma_5 = \begin{pmatrix} 1 & 0 & 0 & 0 \\ 0 & 1 & 0 & 0 \\ 0 & 0 & -1 & 0 \\ 0 & 0 & 0 & -1 \end{pmatrix}.
  \end{gathered}
\end{equation}
In addition to the anti-commutation relation (A.18) the gamma matrices obey
(here $\mu = 1, \ldots, 5$)
\begin{equation}
  \gamma_\mu = \gamma_\mu^\dagger = \gamma_\mu^{-1}.
\end{equation}

When we discuss charge conjugation, we need the charge conjugation matrix $C$
defined through the relations ($\mu = 1, \ldots, 4$)
\begin{equation}
  C\gamma_\mu C^{-1} = -\gamma_\mu^T .
\end{equation}
Using the explicit form of the gamma matrices it is easy to see that in the
chiral representation the charge conjugation matrix is given by
\begin{equation}
  C = i\gamma_2\gamma_4 .
\end{equation}
It obeys
\begin{equation}
  C = C^{-1} = C^\dagger = -C^T .
\end{equation}

We finally quote a simple formula for the inverse of linear combinations of
gamma matrices ($a, b_\mu \in \mathbb{R}$):
\begin{equation}
  \left( a\,\mathbb{1} + i\sum_{\mu=1}^{4}\gamma_\mu b_\mu \right)^{-1}
  = \frac{ a\,\mathbb{1} - i\sum_{\mu=1}^{4}\gamma_\mu b_\mu }{ a^2 + \sum_{\mu=1}^{4}b_\mu^2 }.
\end{equation}
This formula can be verified by multiplying both sides with
$a\mathbb{1} + i\sum_\mu \gamma_\mu b_\mu$.

\section{Fourier transformation on the lattice}

The goal of this appendix is to discuss the Fourier transform $\tilde{f}(p)$
of functions $f(n)$ defined on the lattice $\Lambda$. The lattice is given by
\begin{equation}
  \Lambda = \left\{ n = (n_1, n_2, n_3, n_4) \,\big|\, n_\mu = 0, 1, \ldots, N_\mu - 1 \right\},
\end{equation}
and in most of our applications we have $N_1 = N_2 = N_3 = N$, $N_4 = N_T$.
For the total number of lattice points we introduce the abbreviation
\begin{equation}
  |\Lambda| = N_1 N_2 N_3 N_4 .
\end{equation}

We impose toroidal boundary conditions
\begin{equation}
  f(n + \hat{\mu}N_\mu) = e^{i 2\pi\theta_\mu}\, f(n)
\end{equation}
for each of the directions $\mu$. Here $\hat{\mu}$ denotes the unit vector in
$\mu$-direction. Directions with periodic boundary conditions have
$\theta_\mu = 0$, anti-periodic boundary conditions correspond to
$\theta_\mu = 1/2$.

The momentum space $\tilde{\Lambda}$, which corresponds to the lattice
$\Lambda$ with the boundary conditions (A.29), is defined as
\begin{equation}
  \tilde{\Lambda} = \left\{ p = (p_1, p_2, p_3, p_4) \,\Big|\, p_\mu =
  \frac{2\pi}{aN_\mu}(k_\mu + \theta_\mu),\; k_\mu = -\frac{N_\mu}{2} + 1, \ldots, \frac{N_\mu}{2} \right\}.
\end{equation}
The boundary phases $\theta_\mu$ have to be included in the definition of the
momenta $p_\mu$ such that the plane waves
\begin{equation}
  \exp(i\,p\cdot n\,a) \qquad \text{with} \qquad p\cdot n = \sum_{\mu=1}^{4}p_\mu n_\mu
\end{equation}
also obey the boundary conditions (A.29).

The basic formula, underlying Fourier transformation on the lattice, is (here
$l$ is an integer with $0 \le l \le N - 1$)
\begin{equation}
  \frac{1}{N}\sum_{j=-N/2+1}^{N/2}\exp\!\left(i\,\frac{2\pi}{N}\,lj\right)
  = \frac{1}{N}\sum_{j=0}^{N-1}\exp\!\left(i\,\frac{2\pi}{N}\,lj\right) = \delta_{l0}.
\end{equation}
For $l = 0$ this formula is trivial. For $l \ne 0$ this formula follows from
applying the well-known algebraic identity
\begin{equation}
  \sum_{j=0}^{N-1}q^j = \frac{1 - q^N}{1 - q} \qquad \text{to} \qquad
  q = \exp\!\left(i\,\frac{2\pi}{N}\,l\right).
\end{equation}

We can combine four of the 1D sums in (A.32) to obtain the following
identities:
\begin{equation}
  \frac{1}{|\Lambda|}\sum_{p\in\tilde{\Lambda}}\exp\!\left(i\,p\cdot(n-n')\,a\right)
  = \delta(n-n') = \delta_{n_1n'_1}\delta_{n_2n'_2}\delta_{n_3n'_3}\delta_{n_4n'_4},
\end{equation}
\begin{equation}
  \frac{1}{|\Lambda|}\sum_{n\in\Lambda}\exp\!\left(i\,(p-p')\cdot n\,a\right)
  = \delta(p-p') \equiv \delta_{k_1k'_1}\delta_{k_2k'_2}\delta_{k_3k'_3}\delta_{k_4k'_4}.
\end{equation}
We stress that the right-hand side of (A.35) is a product of four Kronecker
deltas for the integers $k_\mu$ which label the momentum components $p_\mu$
(compare (A.30)).

If we now define the Fourier transform
\begin{equation}
  \tilde{f}(p) = \frac{1}{\sqrt{|\Lambda|}}\sum_{n\in\Lambda}f(n)\,\exp(-i\,p\cdot n\,a),
\end{equation}
we find for the inverse transformation
\begin{equation}
  f(n) = \frac{1}{\sqrt{|\Lambda|}}\sum_{p\in\tilde{\Lambda}}\tilde{f}(p)\,\exp(i\,p\cdot n\,a).
\end{equation}
The last equation follows immediately from inserting the definition of the
Fourier transform in the inverse transformation and using (A.34).

\section{Wilson's formulation of lattice QCD}

In this appendix we collect the defining formulas for Wilson's formulation of
QCD on the lattice. The dynamical variables are the group-valued link
variables $U_\mu(n)$ and the Grassmann-valued fermion fields
$\psi^{(f)}(n)^\alpha_c$, $\bar{\psi}^{(f)}(n)^\alpha_c$. They live on the
links, respectively the sites of our lattice (A.27). Vacuum expectation values
are calculated according to
\begin{equation}
  \ev{O} = \frac{1}{Z}\int D[\psi,\bar{\psi}]\,D[U]\,
  e^{-S_F[\psi,\bar{\psi},U]-S_G[U]}\,O[\psi,\bar{\psi},U],
\end{equation}
where the partition function is given by
\begin{equation}
  Z = \int D[\psi,\bar{\psi}]\,D[U]\,
  e^{-S_F[\psi,\bar{\psi},U]-S_G[U]} .
\end{equation}
The measures over fermion and gauge fields are products over the measures for
the individual field variables:
\begin{align}
  D[\psi,\bar{\psi}] &= \prod_{n\in\Lambda}\prod_{f,\alpha,c}
  d\bar{\psi}^{(f)}(n)^\alpha_c\,d\psi^{(f)}(n)^\alpha_c, \notag\\
  D[U] &= \prod_{n\in\Lambda}\prod_{\mu=1}^{4} dU_\mu(n).
\end{align}
For the individual link variables $U_\mu(n)$ one uses the Haar measure discussed
in Sect.~3.1. For the fermions the rules for Grassmann integration from
Sect.~5.1 apply. The gauge field action for gauge group SU($N$) is given by
\begin{equation}
  S_G[U] = \frac{\beta}{N}\sum_{n\in\Lambda}\sum_{\mu<\nu}\Re\,\tr[1 - U_{\mu\nu}(n)],
\end{equation}
where the plaquettes are defined as
\begin{equation}
  \begin{split}
    U_{\mu\nu}(n) &= U_\mu(n)\,U_\nu(n+\hat{\mu})\,U_{-\mu}(n+\hat{\mu}+\hat{\nu})\,U_{-\nu}(n+\hat{\nu})\\
    &= U_\mu(n)\,U_\nu(n+\hat{\mu})\,U_\mu(n+\hat{\nu})^\dagger\,U_\nu(n)^\dagger .
  \end{split}
\end{equation}
The fermion action is a sum over $N_f$ flavors:
\begin{equation}
  S_F[\psi,\bar{\psi},U] = \sum_{f=1}^{N_f}a^4\sum_{n,m\in\Lambda}
  \bar{\psi}^{(f)}(n)\,D^{(f)}(n|m)\,\psi^{(f)}(m)
\end{equation}
and the lattice Dirac operator is given by
\begin{equation}
  D^{(f)}(n|m)^{\alpha\beta}_{ab} = \left( m^{(f)} + \frac{4}{a} \right)
  \delta_{\alpha\beta}\delta_{ab}\delta_{n,m}
  - \frac{1}{2a}\sum_{\mu=\pm1}^{\pm4}(1-\gamma_\mu)_{\alpha\beta}\,
  U_\mu(n)_{ab}\,\delta_{n+\hat{\mu},m}.
\end{equation}
In the plaquette definition and in the last equation we use the conventions
\begin{equation}
  \gamma_{-\mu} = -\gamma_\mu,\qquad
  U_{-\mu}(n) = U_\mu(n-\hat{\mu})^\dagger,\qquad \mu = 1, 2, 3, 4 .
\end{equation}
We remark that Wilson's Dirac operator is $\gamma_5$-hermitian, i.e., it obeys
\begin{equation}
  \gamma_5 D \gamma_5 = D^\dagger .
\end{equation}

\section{A few formulas for matrix algebra}

In quantum mechanics one usually deals with hermitian or unitary matrices,
while in lattice QCD often more general matrices occur. In this appendix we
list a few results for general complex matrices together with short remarks
concerning their proof (for a more detailed account see, e.g., [5]).

The basic result for general complex matrices is that they are unitarily
equivalent to upper triangular matrices: Let $M$ be a complex-valued
$N\times N$ matrix. Then there exists a unitary matrix $U$ and an upper
triangular matrix $T$, such that
\begin{equation}
  U^\dagger M U = T .
\end{equation}
This result can be proven by induction in $N$. The elements $t_j$ on the
diagonal of $T$ are the roots of the characteristic polynomial of $M$ since
\begin{equation}
  P(\lambda) = \det[M - \lambda\,\mathbb{1}] = \det[T - \lambda\,\mathbb{1}]
  = \prod_{j=1}^{N}(t_j - \lambda).
\end{equation}

An important consequence of this result is a unique classification of matrices
that can be diagonalized with a unitary transformation. A complex matrix $M$
is called normal if it commutes with its hermitian conjugate, i.e.,
$[M, M^\dagger] = 0$. It is obvious that hermitian or unitary matrices are
normal. The announced result is: If and only if $M$ is normal, then there
exists a unitary matrix $U$ such that
\begin{equation}
  U^\dagger M U = D ,
\end{equation}
where $D$ is diagonal. It is straightforward to see that a matrix $M$ which is
unitarily equivalent to a diagonal matrix is normal. To prove the other
direction we first note that the normality of $M$ implies the normality of the
triangular matrix $T$ corresponding to $M$. When evaluating explicitly the two
sides of the normality condition, $T^\dagger T = T T^\dagger$, for the upper
triangular matrix $T$, one concludes that $T$ must be diagonal and the
statement is proven. Equations (A.49) and (A.48) imply that a normal matrix
has a complete set of orthonormal eigenvectors, the columns of $U$.

The existence of a complete orthonormal set of eigenvectors $v^{(j)}$ with
eigenvalues $\lambda^{(j)}$ can be used to represent the matrix $M$ in the
form
\begin{equation}
  M = \sum_{j=1}^{N}\lambda^{(j)}\,v^{(j)}v^{(j)\dagger},
\end{equation}
the so-called spectral representation. On the right-hand side of this equation
matrix/vector notation was used to write the dyadic product
$v^{(j)}v^{(j)\dagger}$. The spectral representation of the matrix can be used
to define a function of $M$ in terms of the function for the eigenvalues, if
this exists. This gives rise to the spectral theorem
\begin{equation}
  f(M) = \sum_{j=1}^{N}f\!\left(\lambda^{(j)}\right) v^{(j)}v^{(j)\dagger}.
\end{equation}

We finally discuss a formula for the expansion of the determinant:
\begin{equation}
  \det[\mathbb{1} - M] = \exp\!\left(\tr[\ln(\mathbb{1} - M)]\right).
\end{equation}
In this equation $M$ is a complex matrix and the logarithm (where it exists)
is defined through its series expansion. The proof of this result applies the
triangularization (A.47):
\begin{align}
  \det[\mathbb{1} - M] &= \det[\mathbb{1} - T] = \prod_{j=1}^{N}(1 - t_j)
  = \exp\!\left(\sum_{j=1}^{N}\ln(1 - t_j)\right) \notag\\
  &= \exp\!\left(-\sum_{j=1}^{N}\sum_{n=1}^{\infty}\frac{1}{n}(t_j)^n\right)
  = \exp\!\left(-\sum_{n=1}^{\infty}\frac{1}{n}\tr[T^n]\right) \notag\\
  &= \exp\!\left(\tr[\ln(\mathbb{1} - M)]\right).
\end{align}
In the fifth step we have used the fact that when evaluating powers of a
triangular matrix the diagonal elements do not mix with other entries of the
matrix. In the last step we used $\tr[T^n] = \tr[M^n]$ which follows from
(A.47). Since a matrix $A$ may always be written as $A = \mathbb{1} - M$, the
result is often stated as
\begin{equation}
  \det[A] = \exp\!\left(\tr[\ln A]\right).
\end{equation}

\section*{References}

\begin{enumerate}[label={[\arabic*]}]
  \item H. Georgi: Lie Algebras in Particle Physics (Benjamin/Cummings,
    Reading, Massachusetts 1982)
  \item H. F. Jones: Groups, Representations and Physics (Hilger, Bristol 1990)
  \item M. Hamermesh: Group Theory and Its Application to Physical Problems
    (Addison-Wesley, Reading, Massachusetts 1964)
  \item M. E. Peskin and D. V. Schroeder: An Introduction to Quantum Field
    Theory (Addison-Wesley, Reading, Massachusetts 1995)
  \item P. Lancaster: Theory of Matrices (Academic Press, New York 1969)
\end{enumerate}
, so the chapter is
%  automatically lettered "A" and the sections become A.1, A.2, ...
% =====================================================================

\chapter{Appendix}

\section{The Lie groups SU($N$)}

In this appendix we collect basic definitions and conventions for the Lie
groups SU($N$) -- the special unitary groups -- and the corresponding Lie
algebras $su(N)$. For a more detailed presentation we refer the reader to
[1--3].

\subsection{Basic properties}

The defining representation of SU($N$) is given by complex $N\times N$
matrices which are unitary and have determinant 1. This set of matrices is
closed under matrix multiplication: Let $\Omega_1$ and $\Omega_2$ be elements
of SU($N$), i.e., they obey $\Omega_i^\dagger = \Omega_i^{-1}$ and
$\det[\Omega_i] = 1$. Using standard linear algebra manipulations we obtain
\begin{equation}
  \begin{gathered}
    (\Omega_1\Omega_2)^\dagger = \Omega_2^\dagger\Omega_1^\dagger
    = \Omega_2^{-1}\Omega_1^{-1} = (\Omega_1\Omega_2)^{-1},\\[2pt]
    \det[\Omega_1\Omega_2] = \det[\Omega_1]\det[\Omega_2] = 1
  \end{gathered}
\end{equation}
and have thus established that also the product of two SU($N$) matrices is an
SU($N$) matrix. The unit matrix is also in SU($N$) and for each matrix in
SU($N$) there exists an inverse (the hermitian conjugate matrix). Thus, the
set SU($N$) forms a group. Since the group operation -- the matrix
multiplication -- is not commutative the groups SU($N$) are non-abelian
groups.

\subsection{Lie algebra}

Let us now count how many real parameters are needed to describe the matrices
in SU($N$). A complex $N\times N$ matrix has $2N^2$ real parameters. The
requirement of unitarity introduces $N^2$ independent conditions which the
parameters have to obey. One more parameter is used for obeying the determinant
condition such that one needs a total of $N^2-1$ real parameters for describing
SU($N$) matrices.

A convenient way of representing SU($N$) matrices is to write them as
exponentials of basis matrices $T_j$, the so-called generators. In particular,
we write an element $\Omega$ of SU($N$) as
\begin{equation}
  \Omega = \exp\!\left(i\sum_{j=1}^{N^2-1}\omega^{(j)} T_j\right),
\end{equation}
where $\omega^{(j)}$, $j = 1, 2, \ldots, N^2-1$, are the real numbers needed to
parameterize $\Omega$. We remark that the parameters $\omega^{(j)}$ can be
changed continuously, making SU($N$) so-called Lie groups that are groups whose
elements depend continuously on their parameters. In order to cover all of the
group space, the parameters have to be varied only over finite intervals,
making the groups SU($N$) so-called compact Lie groups.

The generators $T_j$, $j = 1, 2, \ldots, N^2-1$, are chosen as traceless,
complex, and hermitian $N\times N$ matrices obeying the normalization condition
\begin{equation}
  \tr[T_j T_k] = \half\,\delta_{jk}.
\end{equation}
In addition, they are related among each other by an algebra of commutation
relations
\begin{equation}
  [T_j, T_k] = i\, f_{jkl}\, T_l .
\end{equation}
The completely anti-symmetric coefficients $f_{jkl}$ are the so-called
structure constants. Below we will give an explicit representation of the
generators for the groups SU(2) and SU(3).

Let us verify that the representation of $\Omega$ indeed describes elements of
SU($N$). Using the facts that the generators are hermitian and that the
$\omega^{(j)}$ are real, one finds that hermitian conjugation of the right-hand
side of the exponent simply produces an extra minus sign in the exponent (from
the complex conjugation of $i$). Thus, it is obvious that $\Omega^\dagger =
\Omega^{-1}$. To show that the determinant equals 1, we use the equation
\begin{equation}
  \det\Omega = \exp(\tr[\ln\Omega])
  = \exp\!\left(i\sum_{j=1}^{N^2-1}\omega^{(j)}\tr T_j\right) = e^0 = 1,
\end{equation}
where in first step we have used a formula for the determinant (see (A.54)
below) and in the third step we have used the fact that the $T_j$ are
traceless.

Not only the group elements but also the exponents of our representation have
an interesting structure. The linear combinations
\begin{equation}
  \sum_{j=1}^{N^2-1}\omega^{(j)} T_j
\end{equation}
of the $T_j$ form the so-called Lie algebra $su(N)$. Their commutation
properties are governed by the commutation relations. Elements of $su(N)$ are
also complex $N\times N$ matrices but have properties different from the
elements of the group. One important difference is the fact that the unit
matrix is contained in the group (all $\omega^{(j)} = 0$), while it is not an
element of the algebra (all $T_j$ are traceless).

\subsection{Generators for SU(2) and SU(3)}

The standard representation of the generators for SU(2) is given by
\begin{equation}
  T_j = \half\,\sigma_j,
\end{equation}
with the Pauli matrices
\begin{equation}
  \sigma_1 = \begin{pmatrix} 0 & 1 \\ 1 & 0 \end{pmatrix},\qquad
  \sigma_2 = \begin{pmatrix} 0 & -i \\ i & 0 \end{pmatrix},\qquad
  \sigma_3 = \begin{pmatrix} 1 & 0 \\ 0 & -1 \end{pmatrix}.
\end{equation}
In this case the structure constants are particularly simple, given by the
completely anti-symmetric tensor, i.e., $f_{jkl} = \varepsilon_{jkl}$.

For SU(3) the generators are given by
\begin{equation}
  T_j = \half\,\lambda_j .
\end{equation}
The Gell--Mann matrices $\lambda_j$ are $3\times 3$ generalizations of the
Pauli matrices:
\begin{equation}
  \begin{gathered}
    \lambda_1 = \begin{pmatrix} 0 & 1 & 0 \\ 1 & 0 & 0 \\ 0 & 0 & 0 \end{pmatrix},\qquad
    \lambda_2 = \begin{pmatrix} 0 & -i & 0 \\ i & 0 & 0 \\ 0 & 0 & 0 \end{pmatrix},\qquad
    \lambda_3 = \begin{pmatrix} 1 & 0 & 0 \\ 0 & -1 & 0 \\ 0 & 0 & 0 \end{pmatrix},\\[4pt]
    \lambda_4 = \begin{pmatrix} 0 & 0 & 1 \\ 0 & 0 & 0 \\ 1 & 0 & 0 \end{pmatrix},\qquad
    \lambda_5 = \begin{pmatrix} 0 & 0 & -i \\ 0 & 0 & 0 \\ i & 0 & 0 \end{pmatrix},\qquad
    \lambda_6 = \begin{pmatrix} 0 & 0 & 0 \\ 0 & 0 & 1 \\ 0 & 1 & 0 \end{pmatrix},\\[4pt]
    \lambda_7 = \begin{pmatrix} 0 & 0 & 0 \\ 0 & 0 & -i \\ 0 & i & 0 \end{pmatrix},\qquad
    \lambda_8 = \frac{1}{\sqrt{3}}
    \begin{pmatrix} 1 & 0 & 0 \\ 0 & 1 & 0 \\ 0 & 0 & -2 \end{pmatrix}.
  \end{gathered}
\end{equation}

\subsection{Derivatives of group elements}

Let us now show an important property of derivatives of group elements. If
$\Omega(\omega)$ is an element of SU($N$) then
\begin{equation}
  M_k(\omega) = i\,\frac{\partial\Omega(\omega)}{\partial\omega^{(k)}}\,
  \Omega(\omega)^\dagger \in su(N),
\end{equation}
i.e., the derivative times the hermitian conjugate is in the Lie algebra. In
order to prove this statement we have to show the defining properties of Lie
algebra elements, i.e., we must show that $M_k(\omega)$ is hermitian and
traceless.

Showing the hermiticity of $M_k(\omega)$ is straightforward. By differentiating
$\Omega(\omega)\Omega(\omega)^\dagger = 1$ with respect to $\omega^{(k)}$ one
finds
\begin{equation}
  \frac{\partial\Omega(\omega)}{\partial\omega^{(k)}}\,\Omega(\omega)^\dagger
  + \Omega(\omega)\,\frac{\partial\Omega(\omega)^\dagger}{\partial\omega^{(k)}} = 0.
\end{equation}
Thus
\begin{align}
  M_k(\omega)^\dagger &=
  \left[ i\,\frac{\partial\Omega(\omega)}{\partial\omega^{(k)}}\,
  \Omega(\omega)^\dagger \right]^\dagger
  = -i\,\Omega(\omega)\,\frac{\partial\Omega(\omega)^\dagger}{\partial\omega^{(k)}} \notag\\
  &= i\,\frac{\partial\Omega(\omega)}{\partial\omega^{(k)}}\,
  \Omega(\omega)^\dagger = M_k(\omega),
\end{align}
where we used (A.12) in the third step.

In order to show that $M_k(\omega)$ is traceless we use the fact that the
determinant of an SU($N$) matrix equals 1 and we differentiate
$\det[\Omega(\omega)]$ with respect to $\omega^{(k)}$. We obtain
\begin{align}
  0 &= \frac{\partial\det[\Omega(\omega)]}{\partial\omega^{(k)}}
  = \frac{\partial\det[\Omega(\omega)]}{\partial\Omega(\omega)_{ab}}\,
  \frac{\partial\Omega(\omega)_{ab}}{\partial\omega^{(k)}} \notag\\
  &= \det[\Omega(\omega)]\left(\Omega(\omega)^{-1}\right)_{ba}\,
  \frac{\partial\Omega(\omega)_{ab}}{\partial\omega^{(k)}}
  = \tr\!\left[\frac{\partial\Omega(\omega)}{\partial\omega^{(k)}}\,
  \Omega(\omega)^\dagger\right],
\end{align}
where in the first step we applied the chain rule for derivatives. In the second
step we used a standard formula for the derivative of the determinant
$\det[\Omega]$ with respect to an entry $\Omega_{ab}$ of the matrix $\Omega$.
In the third step we used $\det[\Omega] = 1$ and $\Omega^{-1} = \Omega^\dagger$.
Equation (A.14) establishes that $M_k(\omega)$ is also traceless and thus we
have shown $M_k(\omega)\in su(N)$.

From (A.11) it follows that for the gauge transformation matrices $\Omega(x)$
with coefficients $\omega^{(k)}(x)$, depending on the space--time coordinate
$x$, the combination $i(\partial_\mu\Omega(x))\Omega(x)^\dagger$ is in the Lie
algebra, since
\begin{equation}
  i(\partial_\mu\Omega(x))\,\Omega(x)^\dagger
  = \sum_k i\,\frac{\partial\Omega(\omega(x))}{\partial\omega^{(k)}(x)}\,
  \Omega(\omega(x))^\dagger\,\partial_\mu\omega^{(k)}(x),
\end{equation}
and the right-hand side is a linear combination of $su(N)$ elements with real
coefficients $\partial_\mu\omega^{(k)}(x)$.

\section{Gamma matrices}

The Euclidean gamma matrices $\gamma_\mu$, $\mu = 1, 2, 3, 4$ can be
constructed from the Minkowski gamma matrices $\gamma_\mu^M$, $\mu = 0, 1, 2,
3$. The latter obey
\begin{equation}
  \{\gamma_\mu^M, \gamma_\nu^M\} = 2\,g_{\mu\nu}\,\mathbb{1},
\end{equation}
with the metric tensor given by $g_{\mu\nu} = \mathrm{diag}(1,-1,-1,-1)$ and
$\mathbb{1}$ is the $4\times 4$ unit matrix. Thus when we define the Euclidean
matrices $\gamma_\mu$ by setting
\begin{equation}
  \gamma_1 = -i\gamma_1^M,\quad
  \gamma_2 = -i\gamma_2^M,\quad
  \gamma_3 = -i\gamma_3^M,\quad
  \gamma_4 = \gamma_0^M,
\end{equation}
we obtain the Euclidean anti-commutation relations
\begin{equation}
  \{\gamma_\mu, \gamma_\nu\} = 2\,\delta_{\mu\nu}\,\mathbb{1}.
\end{equation}

In addition to the matrices $\gamma_\mu$, $\mu = 1, 2, 3, 4$ we define the
matrix $\gamma_5$ as the product
\begin{equation}
  \gamma_5 = \gamma_1\gamma_2\gamma_3\gamma_4 .
\end{equation}
The matrix $\gamma_5$ anti-commutes with all other gamma matrices $\gamma_\mu$,
$\mu = 1, 2, 3, 4$ and obeys $\gamma_5^2 = 1$.

An explicit representation of the Euclidean gamma matrices can be obtained from
a representation of the Minkowski gamma matrices (see, e.g., [4]). Here we give
the so-called chiral representation where $\gamma_5$ (the chirality operator)
is diagonal:
\begin{equation}
  \gamma_{1,2,3} = \begin{pmatrix} 0 & -i\sigma_{1,2,3} \\ i\sigma_{1,2,3} & 0 \end{pmatrix},\qquad
  \gamma_4 = \begin{pmatrix} 0 & \mathbb{1}_2 \\ \mathbb{1}_2 & 0 \end{pmatrix},\qquad
  \gamma_5 = \begin{pmatrix} \mathbb{1}_2 & 0 \\ 0 & -\mathbb{1}_2 \end{pmatrix},
\end{equation}
where the $\sigma_j$ are the Pauli matrices and $\mathbb{1}_2$ is the
$2\times 2$ unit matrix. More explicitly the Euclidean gamma matrices read
\begin{equation}
  \begin{gathered}
    \gamma_1 = \begin{pmatrix} 0 & 0 & 0 & -i \\ 0 & 0 & -i & 0 \\ 0 & i & 0 & 0 \\ i & 0 & 0 & 0 \end{pmatrix},\qquad
    \gamma_2 = \begin{pmatrix} 0 & 0 & 0 & -1 \\ 0 & 0 & 1 & 0 \\ 0 & 1 & 0 & 0 \\ -1 & 0 & 0 & 0 \end{pmatrix},\\[4pt]
    \gamma_3 = \begin{pmatrix} 0 & 0 & -i & 0 \\ 0 & 0 & 0 & i \\ i & 0 & 0 & 0 \\ 0 & -i & 0 & 0 \end{pmatrix},\qquad
    \gamma_4 = \begin{pmatrix} 0 & 0 & 1 & 0 \\ 0 & 0 & 0 & 1 \\ 1 & 0 & 0 & 0 \\ 0 & 1 & 0 & 0 \end{pmatrix},\\[4pt]
    \gamma_5 = \begin{pmatrix} 1 & 0 & 0 & 0 \\ 0 & 1 & 0 & 0 \\ 0 & 0 & -1 & 0 \\ 0 & 0 & 0 & -1 \end{pmatrix}.
  \end{gathered}
\end{equation}
In addition to the anti-commutation relation (A.18) the gamma matrices obey
(here $\mu = 1, \ldots, 5$)
\begin{equation}
  \gamma_\mu = \gamma_\mu^\dagger = \gamma_\mu^{-1}.
\end{equation}

When we discuss charge conjugation, we need the charge conjugation matrix $C$
defined through the relations ($\mu = 1, \ldots, 4$)
\begin{equation}
  C\gamma_\mu C^{-1} = -\gamma_\mu^T .
\end{equation}
Using the explicit form of the gamma matrices it is easy to see that in the
chiral representation the charge conjugation matrix is given by
\begin{equation}
  C = i\gamma_2\gamma_4 .
\end{equation}
It obeys
\begin{equation}
  C = C^{-1} = C^\dagger = -C^T .
\end{equation}

We finally quote a simple formula for the inverse of linear combinations of
gamma matrices ($a, b_\mu \in \mathbb{R}$):
\begin{equation}
  \left( a\,\mathbb{1} + i\sum_{\mu=1}^{4}\gamma_\mu b_\mu \right)^{-1}
  = \frac{ a\,\mathbb{1} - i\sum_{\mu=1}^{4}\gamma_\mu b_\mu }{ a^2 + \sum_{\mu=1}^{4}b_\mu^2 }.
\end{equation}
This formula can be verified by multiplying both sides with
$a\mathbb{1} + i\sum_\mu \gamma_\mu b_\mu$.

\section{Fourier transformation on the lattice}

The goal of this appendix is to discuss the Fourier transform $\tilde{f}(p)$
of functions $f(n)$ defined on the lattice $\Lambda$. The lattice is given by
\begin{equation}
  \Lambda = \left\{ n = (n_1, n_2, n_3, n_4) \,\big|\, n_\mu = 0, 1, \ldots, N_\mu - 1 \right\},
\end{equation}
and in most of our applications we have $N_1 = N_2 = N_3 = N$, $N_4 = N_T$.
For the total number of lattice points we introduce the abbreviation
\begin{equation}
  |\Lambda| = N_1 N_2 N_3 N_4 .
\end{equation}

We impose toroidal boundary conditions
\begin{equation}
  f(n + \hat{\mu}N_\mu) = e^{i 2\pi\theta_\mu}\, f(n)
\end{equation}
for each of the directions $\mu$. Here $\hat{\mu}$ denotes the unit vector in
$\mu$-direction. Directions with periodic boundary conditions have
$\theta_\mu = 0$, anti-periodic boundary conditions correspond to
$\theta_\mu = 1/2$.

The momentum space $\tilde{\Lambda}$, which corresponds to the lattice
$\Lambda$ with the boundary conditions (A.29), is defined as
\begin{equation}
  \tilde{\Lambda} = \left\{ p = (p_1, p_2, p_3, p_4) \,\Big|\, p_\mu =
  \frac{2\pi}{aN_\mu}(k_\mu + \theta_\mu),\; k_\mu = -\frac{N_\mu}{2} + 1, \ldots, \frac{N_\mu}{2} \right\}.
\end{equation}
The boundary phases $\theta_\mu$ have to be included in the definition of the
momenta $p_\mu$ such that the plane waves
\begin{equation}
  \exp(i\,p\cdot n\,a) \qquad \text{with} \qquad p\cdot n = \sum_{\mu=1}^{4}p_\mu n_\mu
\end{equation}
also obey the boundary conditions (A.29).

The basic formula, underlying Fourier transformation on the lattice, is (here
$l$ is an integer with $0 \le l \le N - 1$)
\begin{equation}
  \frac{1}{N}\sum_{j=-N/2+1}^{N/2}\exp\!\left(i\,\frac{2\pi}{N}\,lj\right)
  = \frac{1}{N}\sum_{j=0}^{N-1}\exp\!\left(i\,\frac{2\pi}{N}\,lj\right) = \delta_{l0}.
\end{equation}
For $l = 0$ this formula is trivial. For $l \ne 0$ this formula follows from
applying the well-known algebraic identity
\begin{equation}
  \sum_{j=0}^{N-1}q^j = \frac{1 - q^N}{1 - q} \qquad \text{to} \qquad
  q = \exp\!\left(i\,\frac{2\pi}{N}\,l\right).
\end{equation}

We can combine four of the 1D sums in (A.32) to obtain the following
identities:
\begin{equation}
  \frac{1}{|\Lambda|}\sum_{p\in\tilde{\Lambda}}\exp\!\left(i\,p\cdot(n-n')\,a\right)
  = \delta(n-n') = \delta_{n_1n'_1}\delta_{n_2n'_2}\delta_{n_3n'_3}\delta_{n_4n'_4},
\end{equation}
\begin{equation}
  \frac{1}{|\Lambda|}\sum_{n\in\Lambda}\exp\!\left(i\,(p-p')\cdot n\,a\right)
  = \delta(p-p') \equiv \delta_{k_1k'_1}\delta_{k_2k'_2}\delta_{k_3k'_3}\delta_{k_4k'_4}.
\end{equation}
We stress that the right-hand side of (A.35) is a product of four Kronecker
deltas for the integers $k_\mu$ which label the momentum components $p_\mu$
(compare (A.30)).

If we now define the Fourier transform
\begin{equation}
  \tilde{f}(p) = \frac{1}{\sqrt{|\Lambda|}}\sum_{n\in\Lambda}f(n)\,\exp(-i\,p\cdot n\,a),
\end{equation}
we find for the inverse transformation
\begin{equation}
  f(n) = \frac{1}{\sqrt{|\Lambda|}}\sum_{p\in\tilde{\Lambda}}\tilde{f}(p)\,\exp(i\,p\cdot n\,a).
\end{equation}
The last equation follows immediately from inserting the definition of the
Fourier transform in the inverse transformation and using (A.34).

\section{Wilson's formulation of lattice QCD}

In this appendix we collect the defining formulas for Wilson's formulation of
QCD on the lattice. The dynamical variables are the group-valued link
variables $U_\mu(n)$ and the Grassmann-valued fermion fields
$\psi^{(f)}(n)^\alpha_c$, $\bar{\psi}^{(f)}(n)^\alpha_c$. They live on the
links, respectively the sites of our lattice (A.27). Vacuum expectation values
are calculated according to
\begin{equation}
  \ev{O} = \frac{1}{Z}\int D[\psi,\bar{\psi}]\,D[U]\,
  e^{-S_F[\psi,\bar{\psi},U]-S_G[U]}\,O[\psi,\bar{\psi},U],
\end{equation}
where the partition function is given by
\begin{equation}
  Z = \int D[\psi,\bar{\psi}]\,D[U]\,
  e^{-S_F[\psi,\bar{\psi},U]-S_G[U]} .
\end{equation}
The measures over fermion and gauge fields are products over the measures for
the individual field variables:
\begin{align}
  D[\psi,\bar{\psi}] &= \prod_{n\in\Lambda}\prod_{f,\alpha,c}
  d\bar{\psi}^{(f)}(n)^\alpha_c\,d\psi^{(f)}(n)^\alpha_c, \notag\\
  D[U] &= \prod_{n\in\Lambda}\prod_{\mu=1}^{4} dU_\mu(n).
\end{align}
For the individual link variables $U_\mu(n)$ one uses the Haar measure discussed
in Sect.~3.1. For the fermions the rules for Grassmann integration from
Sect.~5.1 apply. The gauge field action for gauge group SU($N$) is given by
\begin{equation}
  S_G[U] = \frac{\beta}{N}\sum_{n\in\Lambda}\sum_{\mu<\nu}\Re\,\tr[1 - U_{\mu\nu}(n)],
\end{equation}
where the plaquettes are defined as
\begin{equation}
  \begin{split}
    U_{\mu\nu}(n) &= U_\mu(n)\,U_\nu(n+\hat{\mu})\,U_{-\mu}(n+\hat{\mu}+\hat{\nu})\,U_{-\nu}(n+\hat{\nu})\\
    &= U_\mu(n)\,U_\nu(n+\hat{\mu})\,U_\mu(n+\hat{\nu})^\dagger\,U_\nu(n)^\dagger .
  \end{split}
\end{equation}
The fermion action is a sum over $N_f$ flavors:
\begin{equation}
  S_F[\psi,\bar{\psi},U] = \sum_{f=1}^{N_f}a^4\sum_{n,m\in\Lambda}
  \bar{\psi}^{(f)}(n)\,D^{(f)}(n|m)\,\psi^{(f)}(m)
\end{equation}
and the lattice Dirac operator is given by
\begin{equation}
  D^{(f)}(n|m)^{\alpha\beta}_{ab} = \left( m^{(f)} + \frac{4}{a} \right)
  \delta_{\alpha\beta}\delta_{ab}\delta_{n,m}
  - \frac{1}{2a}\sum_{\mu=\pm1}^{\pm4}(1-\gamma_\mu)_{\alpha\beta}\,
  U_\mu(n)_{ab}\,\delta_{n+\hat{\mu},m}.
\end{equation}
In the plaquette definition and in the last equation we use the conventions
\begin{equation}
  \gamma_{-\mu} = -\gamma_\mu,\qquad
  U_{-\mu}(n) = U_\mu(n-\hat{\mu})^\dagger,\qquad \mu = 1, 2, 3, 4 .
\end{equation}
We remark that Wilson's Dirac operator is $\gamma_5$-hermitian, i.e., it obeys
\begin{equation}
  \gamma_5 D \gamma_5 = D^\dagger .
\end{equation}

\section{A few formulas for matrix algebra}

In quantum mechanics one usually deals with hermitian or unitary matrices,
while in lattice QCD often more general matrices occur. In this appendix we
list a few results for general complex matrices together with short remarks
concerning their proof (for a more detailed account see, e.g., [5]).

The basic result for general complex matrices is that they are unitarily
equivalent to upper triangular matrices: Let $M$ be a complex-valued
$N\times N$ matrix. Then there exists a unitary matrix $U$ and an upper
triangular matrix $T$, such that
\begin{equation}
  U^\dagger M U = T .
\end{equation}
This result can be proven by induction in $N$. The elements $t_j$ on the
diagonal of $T$ are the roots of the characteristic polynomial of $M$ since
\begin{equation}
  P(\lambda) = \det[M - \lambda\,\mathbb{1}] = \det[T - \lambda\,\mathbb{1}]
  = \prod_{j=1}^{N}(t_j - \lambda).
\end{equation}

An important consequence of this result is a unique classification of matrices
that can be diagonalized with a unitary transformation. A complex matrix $M$
is called normal if it commutes with its hermitian conjugate, i.e.,
$[M, M^\dagger] = 0$. It is obvious that hermitian or unitary matrices are
normal. The announced result is: If and only if $M$ is normal, then there
exists a unitary matrix $U$ such that
\begin{equation}
  U^\dagger M U = D ,
\end{equation}
where $D$ is diagonal. It is straightforward to see that a matrix $M$ which is
unitarily equivalent to a diagonal matrix is normal. To prove the other
direction we first note that the normality of $M$ implies the normality of the
triangular matrix $T$ corresponding to $M$. When evaluating explicitly the two
sides of the normality condition, $T^\dagger T = T T^\dagger$, for the upper
triangular matrix $T$, one concludes that $T$ must be diagonal and the
statement is proven. Equations (A.49) and (A.48) imply that a normal matrix
has a complete set of orthonormal eigenvectors, the columns of $U$.

The existence of a complete orthonormal set of eigenvectors $v^{(j)}$ with
eigenvalues $\lambda^{(j)}$ can be used to represent the matrix $M$ in the
form
\begin{equation}
  M = \sum_{j=1}^{N}\lambda^{(j)}\,v^{(j)}v^{(j)\dagger},
\end{equation}
the so-called spectral representation. On the right-hand side of this equation
matrix/vector notation was used to write the dyadic product
$v^{(j)}v^{(j)\dagger}$. The spectral representation of the matrix can be used
to define a function of $M$ in terms of the function for the eigenvalues, if
this exists. This gives rise to the spectral theorem
\begin{equation}
  f(M) = \sum_{j=1}^{N}f\!\left(\lambda^{(j)}\right) v^{(j)}v^{(j)\dagger}.
\end{equation}

We finally discuss a formula for the expansion of the determinant:
\begin{equation}
  \det[\mathbb{1} - M] = \exp\!\left(\tr[\ln(\mathbb{1} - M)]\right).
\end{equation}
In this equation $M$ is a complex matrix and the logarithm (where it exists)
is defined through its series expansion. The proof of this result applies the
triangularization (A.47):
\begin{align}
  \det[\mathbb{1} - M] &= \det[\mathbb{1} - T] = \prod_{j=1}^{N}(1 - t_j)
  = \exp\!\left(\sum_{j=1}^{N}\ln(1 - t_j)\right) \notag\\
  &= \exp\!\left(-\sum_{j=1}^{N}\sum_{n=1}^{\infty}\frac{1}{n}(t_j)^n\right)
  = \exp\!\left(-\sum_{n=1}^{\infty}\frac{1}{n}\tr[T^n]\right) \notag\\
  &= \exp\!\left(\tr[\ln(\mathbb{1} - M)]\right).
\end{align}
In the fifth step we have used the fact that when evaluating powers of a
triangular matrix the diagonal elements do not mix with other entries of the
matrix. In the last step we used $\tr[T^n] = \tr[M^n]$ which follows from
(A.47). Since a matrix $A$ may always be written as $A = \mathbb{1} - M$, the
result is often stated as
\begin{equation}
  \det[A] = \exp\!\left(\tr[\ln A]\right).
\end{equation}

\section*{References}

\begin{enumerate}[label={[\arabic*]}]
  \item H. Georgi: Lie Algebras in Particle Physics (Benjamin/Cummings,
    Reading, Massachusetts 1982)
  \item H. F. Jones: Groups, Representations and Physics (Hilger, Bristol 1990)
  \item M. Hamermesh: Group Theory and Its Application to Physical Problems
    (Addison-Wesley, Reading, Massachusetts 1964)
  \item M. E. Peskin and D. V. Schroeder: An Introduction to Quantum Field
    Theory (Addison-Wesley, Reading, Massachusetts 1995)
  \item P. Lancaster: Theory of Matrices (Academic Press, New York 1969)
\end{enumerate}
, so the chapter is
%  automatically lettered "A" and the sections become A.1, A.2, ...
% =====================================================================

\chapter{Appendix}

\section{The Lie groups SU($N$)}

In this appendix we collect basic definitions and conventions for the Lie
groups SU($N$) -- the special unitary groups -- and the corresponding Lie
algebras $su(N)$. For a more detailed presentation we refer the reader to
[1--3].

\subsection{Basic properties}

The defining representation of SU($N$) is given by complex $N\times N$
matrices which are unitary and have determinant 1. This set of matrices is
closed under matrix multiplication: Let $\Omega_1$ and $\Omega_2$ be elements
of SU($N$), i.e., they obey $\Omega_i^\dagger = \Omega_i^{-1}$ and
$\det[\Omega_i] = 1$. Using standard linear algebra manipulations we obtain
\begin{equation}
  \begin{gathered}
    (\Omega_1\Omega_2)^\dagger = \Omega_2^\dagger\Omega_1^\dagger
    = \Omega_2^{-1}\Omega_1^{-1} = (\Omega_1\Omega_2)^{-1},\\[2pt]
    \det[\Omega_1\Omega_2] = \det[\Omega_1]\det[\Omega_2] = 1
  \end{gathered}
\end{equation}
and have thus established that also the product of two SU($N$) matrices is an
SU($N$) matrix. The unit matrix is also in SU($N$) and for each matrix in
SU($N$) there exists an inverse (the hermitian conjugate matrix). Thus, the
set SU($N$) forms a group. Since the group operation -- the matrix
multiplication -- is not commutative the groups SU($N$) are non-abelian
groups.

\subsection{Lie algebra}

Let us now count how many real parameters are needed to describe the matrices
in SU($N$). A complex $N\times N$ matrix has $2N^2$ real parameters. The
requirement of unitarity introduces $N^2$ independent conditions which the
parameters have to obey. One more parameter is used for obeying the determinant
condition such that one needs a total of $N^2-1$ real parameters for describing
SU($N$) matrices.

A convenient way of representing SU($N$) matrices is to write them as
exponentials of basis matrices $T_j$, the so-called generators. In particular,
we write an element $\Omega$ of SU($N$) as
\begin{equation}
  \Omega = \exp\!\left(i\sum_{j=1}^{N^2-1}\omega^{(j)} T_j\right),
\end{equation}
where $\omega^{(j)}$, $j = 1, 2, \ldots, N^2-1$, are the real numbers needed to
parameterize $\Omega$. We remark that the parameters $\omega^{(j)}$ can be
changed continuously, making SU($N$) so-called Lie groups that are groups whose
elements depend continuously on their parameters. In order to cover all of the
group space, the parameters have to be varied only over finite intervals,
making the groups SU($N$) so-called compact Lie groups.

The generators $T_j$, $j = 1, 2, \ldots, N^2-1$, are chosen as traceless,
complex, and hermitian $N\times N$ matrices obeying the normalization condition
\begin{equation}
  \tr[T_j T_k] = \half\,\delta_{jk}.
\end{equation}
In addition, they are related among each other by an algebra of commutation
relations
\begin{equation}
  [T_j, T_k] = i\, f_{jkl}\, T_l .
\end{equation}
The completely anti-symmetric coefficients $f_{jkl}$ are the so-called
structure constants. Below we will give an explicit representation of the
generators for the groups SU(2) and SU(3).

Let us verify that the representation of $\Omega$ indeed describes elements of
SU($N$). Using the facts that the generators are hermitian and that the
$\omega^{(j)}$ are real, one finds that hermitian conjugation of the right-hand
side of the exponent simply produces an extra minus sign in the exponent (from
the complex conjugation of $i$). Thus, it is obvious that $\Omega^\dagger =
\Omega^{-1}$. To show that the determinant equals 1, we use the equation
\begin{equation}
  \det\Omega = \exp(\tr[\ln\Omega])
  = \exp\!\left(i\sum_{j=1}^{N^2-1}\omega^{(j)}\tr T_j\right) = e^0 = 1,
\end{equation}
where in first step we have used a formula for the determinant (see (A.54)
below) and in the third step we have used the fact that the $T_j$ are
traceless.

Not only the group elements but also the exponents of our representation have
an interesting structure. The linear combinations
\begin{equation}
  \sum_{j=1}^{N^2-1}\omega^{(j)} T_j
\end{equation}
of the $T_j$ form the so-called Lie algebra $su(N)$. Their commutation
properties are governed by the commutation relations. Elements of $su(N)$ are
also complex $N\times N$ matrices but have properties different from the
elements of the group. One important difference is the fact that the unit
matrix is contained in the group (all $\omega^{(j)} = 0$), while it is not an
element of the algebra (all $T_j$ are traceless).

\subsection{Generators for SU(2) and SU(3)}

The standard representation of the generators for SU(2) is given by
\begin{equation}
  T_j = \half\,\sigma_j,
\end{equation}
with the Pauli matrices
\begin{equation}
  \sigma_1 = \begin{pmatrix} 0 & 1 \\ 1 & 0 \end{pmatrix},\qquad
  \sigma_2 = \begin{pmatrix} 0 & -i \\ i & 0 \end{pmatrix},\qquad
  \sigma_3 = \begin{pmatrix} 1 & 0 \\ 0 & -1 \end{pmatrix}.
\end{equation}
In this case the structure constants are particularly simple, given by the
completely anti-symmetric tensor, i.e., $f_{jkl} = \varepsilon_{jkl}$.

For SU(3) the generators are given by
\begin{equation}
  T_j = \half\,\lambda_j .
\end{equation}
The Gell--Mann matrices $\lambda_j$ are $3\times 3$ generalizations of the
Pauli matrices:
\begin{equation}
  \begin{gathered}
    \lambda_1 = \begin{pmatrix} 0 & 1 & 0 \\ 1 & 0 & 0 \\ 0 & 0 & 0 \end{pmatrix},\qquad
    \lambda_2 = \begin{pmatrix} 0 & -i & 0 \\ i & 0 & 0 \\ 0 & 0 & 0 \end{pmatrix},\qquad
    \lambda_3 = \begin{pmatrix} 1 & 0 & 0 \\ 0 & -1 & 0 \\ 0 & 0 & 0 \end{pmatrix},\\[4pt]
    \lambda_4 = \begin{pmatrix} 0 & 0 & 1 \\ 0 & 0 & 0 \\ 1 & 0 & 0 \end{pmatrix},\qquad
    \lambda_5 = \begin{pmatrix} 0 & 0 & -i \\ 0 & 0 & 0 \\ i & 0 & 0 \end{pmatrix},\qquad
    \lambda_6 = \begin{pmatrix} 0 & 0 & 0 \\ 0 & 0 & 1 \\ 0 & 1 & 0 \end{pmatrix},\\[4pt]
    \lambda_7 = \begin{pmatrix} 0 & 0 & 0 \\ 0 & 0 & -i \\ 0 & i & 0 \end{pmatrix},\qquad
    \lambda_8 = \frac{1}{\sqrt{3}}
    \begin{pmatrix} 1 & 0 & 0 \\ 0 & 1 & 0 \\ 0 & 0 & -2 \end{pmatrix}.
  \end{gathered}
\end{equation}

\subsection{Derivatives of group elements}

Let us now show an important property of derivatives of group elements. If
$\Omega(\omega)$ is an element of SU($N$) then
\begin{equation}
  M_k(\omega) = i\,\frac{\partial\Omega(\omega)}{\partial\omega^{(k)}}\,
  \Omega(\omega)^\dagger \in su(N),
\end{equation}
i.e., the derivative times the hermitian conjugate is in the Lie algebra. In
order to prove this statement we have to show the defining properties of Lie
algebra elements, i.e., we must show that $M_k(\omega)$ is hermitian and
traceless.

Showing the hermiticity of $M_k(\omega)$ is straightforward. By differentiating
$\Omega(\omega)\Omega(\omega)^\dagger = 1$ with respect to $\omega^{(k)}$ one
finds
\begin{equation}
  \frac{\partial\Omega(\omega)}{\partial\omega^{(k)}}\,\Omega(\omega)^\dagger
  + \Omega(\omega)\,\frac{\partial\Omega(\omega)^\dagger}{\partial\omega^{(k)}} = 0.
\end{equation}
Thus
\begin{align}
  M_k(\omega)^\dagger &=
  \left[ i\,\frac{\partial\Omega(\omega)}{\partial\omega^{(k)}}\,
  \Omega(\omega)^\dagger \right]^\dagger
  = -i\,\Omega(\omega)\,\frac{\partial\Omega(\omega)^\dagger}{\partial\omega^{(k)}} \notag\\
  &= i\,\frac{\partial\Omega(\omega)}{\partial\omega^{(k)}}\,
  \Omega(\omega)^\dagger = M_k(\omega),
\end{align}
where we used (A.12) in the third step.

In order to show that $M_k(\omega)$ is traceless we use the fact that the
determinant of an SU($N$) matrix equals 1 and we differentiate
$\det[\Omega(\omega)]$ with respect to $\omega^{(k)}$. We obtain
\begin{align}
  0 &= \frac{\partial\det[\Omega(\omega)]}{\partial\omega^{(k)}}
  = \frac{\partial\det[\Omega(\omega)]}{\partial\Omega(\omega)_{ab}}\,
  \frac{\partial\Omega(\omega)_{ab}}{\partial\omega^{(k)}} \notag\\
  &= \det[\Omega(\omega)]\left(\Omega(\omega)^{-1}\right)_{ba}\,
  \frac{\partial\Omega(\omega)_{ab}}{\partial\omega^{(k)}}
  = \tr\!\left[\frac{\partial\Omega(\omega)}{\partial\omega^{(k)}}\,
  \Omega(\omega)^\dagger\right],
\end{align}
where in the first step we applied the chain rule for derivatives. In the second
step we used a standard formula for the derivative of the determinant
$\det[\Omega]$ with respect to an entry $\Omega_{ab}$ of the matrix $\Omega$.
In the third step we used $\det[\Omega] = 1$ and $\Omega^{-1} = \Omega^\dagger$.
Equation (A.14) establishes that $M_k(\omega)$ is also traceless and thus we
have shown $M_k(\omega)\in su(N)$.

From (A.11) it follows that for the gauge transformation matrices $\Omega(x)$
with coefficients $\omega^{(k)}(x)$, depending on the space--time coordinate
$x$, the combination $i(\partial_\mu\Omega(x))\Omega(x)^\dagger$ is in the Lie
algebra, since
\begin{equation}
  i(\partial_\mu\Omega(x))\,\Omega(x)^\dagger
  = \sum_k i\,\frac{\partial\Omega(\omega(x))}{\partial\omega^{(k)}(x)}\,
  \Omega(\omega(x))^\dagger\,\partial_\mu\omega^{(k)}(x),
\end{equation}
and the right-hand side is a linear combination of $su(N)$ elements with real
coefficients $\partial_\mu\omega^{(k)}(x)$.

\section{Gamma matrices}

The Euclidean gamma matrices $\gamma_\mu$, $\mu = 1, 2, 3, 4$ can be
constructed from the Minkowski gamma matrices $\gamma_\mu^M$, $\mu = 0, 1, 2,
3$. The latter obey
\begin{equation}
  \{\gamma_\mu^M, \gamma_\nu^M\} = 2\,g_{\mu\nu}\,\mathbb{1},
\end{equation}
with the metric tensor given by $g_{\mu\nu} = \mathrm{diag}(1,-1,-1,-1)$ and
$\mathbb{1}$ is the $4\times 4$ unit matrix. Thus when we define the Euclidean
matrices $\gamma_\mu$ by setting
\begin{equation}
  \gamma_1 = -i\gamma_1^M,\quad
  \gamma_2 = -i\gamma_2^M,\quad
  \gamma_3 = -i\gamma_3^M,\quad
  \gamma_4 = \gamma_0^M,
\end{equation}
we obtain the Euclidean anti-commutation relations
\begin{equation}
  \{\gamma_\mu, \gamma_\nu\} = 2\,\delta_{\mu\nu}\,\mathbb{1}.
\end{equation}

In addition to the matrices $\gamma_\mu$, $\mu = 1, 2, 3, 4$ we define the
matrix $\gamma_5$ as the product
\begin{equation}
  \gamma_5 = \gamma_1\gamma_2\gamma_3\gamma_4 .
\end{equation}
The matrix $\gamma_5$ anti-commutes with all other gamma matrices $\gamma_\mu$,
$\mu = 1, 2, 3, 4$ and obeys $\gamma_5^2 = 1$.

An explicit representation of the Euclidean gamma matrices can be obtained from
a representation of the Minkowski gamma matrices (see, e.g., [4]). Here we give
the so-called chiral representation where $\gamma_5$ (the chirality operator)
is diagonal:
\begin{equation}
  \gamma_{1,2,3} = \begin{pmatrix} 0 & -i\sigma_{1,2,3} \\ i\sigma_{1,2,3} & 0 \end{pmatrix},\qquad
  \gamma_4 = \begin{pmatrix} 0 & \mathbb{1}_2 \\ \mathbb{1}_2 & 0 \end{pmatrix},\qquad
  \gamma_5 = \begin{pmatrix} \mathbb{1}_2 & 0 \\ 0 & -\mathbb{1}_2 \end{pmatrix},
\end{equation}
where the $\sigma_j$ are the Pauli matrices and $\mathbb{1}_2$ is the
$2\times 2$ unit matrix. More explicitly the Euclidean gamma matrices read
\begin{equation}
  \begin{gathered}
    \gamma_1 = \begin{pmatrix} 0 & 0 & 0 & -i \\ 0 & 0 & -i & 0 \\ 0 & i & 0 & 0 \\ i & 0 & 0 & 0 \end{pmatrix},\qquad
    \gamma_2 = \begin{pmatrix} 0 & 0 & 0 & -1 \\ 0 & 0 & 1 & 0 \\ 0 & 1 & 0 & 0 \\ -1 & 0 & 0 & 0 \end{pmatrix},\\[4pt]
    \gamma_3 = \begin{pmatrix} 0 & 0 & -i & 0 \\ 0 & 0 & 0 & i \\ i & 0 & 0 & 0 \\ 0 & -i & 0 & 0 \end{pmatrix},\qquad
    \gamma_4 = \begin{pmatrix} 0 & 0 & 1 & 0 \\ 0 & 0 & 0 & 1 \\ 1 & 0 & 0 & 0 \\ 0 & 1 & 0 & 0 \end{pmatrix},\\[4pt]
    \gamma_5 = \begin{pmatrix} 1 & 0 & 0 & 0 \\ 0 & 1 & 0 & 0 \\ 0 & 0 & -1 & 0 \\ 0 & 0 & 0 & -1 \end{pmatrix}.
  \end{gathered}
\end{equation}
In addition to the anti-commutation relation (A.18) the gamma matrices obey
(here $\mu = 1, \ldots, 5$)
\begin{equation}
  \gamma_\mu = \gamma_\mu^\dagger = \gamma_\mu^{-1}.
\end{equation}

When we discuss charge conjugation, we need the charge conjugation matrix $C$
defined through the relations ($\mu = 1, \ldots, 4$)
\begin{equation}
  C\gamma_\mu C^{-1} = -\gamma_\mu^T .
\end{equation}
Using the explicit form of the gamma matrices it is easy to see that in the
chiral representation the charge conjugation matrix is given by
\begin{equation}
  C = i\gamma_2\gamma_4 .
\end{equation}
It obeys
\begin{equation}
  C = C^{-1} = C^\dagger = -C^T .
\end{equation}

We finally quote a simple formula for the inverse of linear combinations of
gamma matrices ($a, b_\mu \in \mathbb{R}$):
\begin{equation}
  \left( a\,\mathbb{1} + i\sum_{\mu=1}^{4}\gamma_\mu b_\mu \right)^{-1}
  = \frac{ a\,\mathbb{1} - i\sum_{\mu=1}^{4}\gamma_\mu b_\mu }{ a^2 + \sum_{\mu=1}^{4}b_\mu^2 }.
\end{equation}
This formula can be verified by multiplying both sides with
$a\mathbb{1} + i\sum_\mu \gamma_\mu b_\mu$.

\section{Fourier transformation on the lattice}

The goal of this appendix is to discuss the Fourier transform $\tilde{f}(p)$
of functions $f(n)$ defined on the lattice $\Lambda$. The lattice is given by
\begin{equation}
  \Lambda = \left\{ n = (n_1, n_2, n_3, n_4) \,\big|\, n_\mu = 0, 1, \ldots, N_\mu - 1 \right\},
\end{equation}
and in most of our applications we have $N_1 = N_2 = N_3 = N$, $N_4 = N_T$.
For the total number of lattice points we introduce the abbreviation
\begin{equation}
  |\Lambda| = N_1 N_2 N_3 N_4 .
\end{equation}

We impose toroidal boundary conditions
\begin{equation}
  f(n + \hat{\mu}N_\mu) = e^{i 2\pi\theta_\mu}\, f(n)
\end{equation}
for each of the directions $\mu$. Here $\hat{\mu}$ denotes the unit vector in
$\mu$-direction. Directions with periodic boundary conditions have
$\theta_\mu = 0$, anti-periodic boundary conditions correspond to
$\theta_\mu = 1/2$.

The momentum space $\tilde{\Lambda}$, which corresponds to the lattice
$\Lambda$ with the boundary conditions (A.29), is defined as
\begin{equation}
  \tilde{\Lambda} = \left\{ p = (p_1, p_2, p_3, p_4) \,\Big|\, p_\mu =
  \frac{2\pi}{aN_\mu}(k_\mu + \theta_\mu),\; k_\mu = -\frac{N_\mu}{2} + 1, \ldots, \frac{N_\mu}{2} \right\}.
\end{equation}
The boundary phases $\theta_\mu$ have to be included in the definition of the
momenta $p_\mu$ such that the plane waves
\begin{equation}
  \exp(i\,p\cdot n\,a) \qquad \text{with} \qquad p\cdot n = \sum_{\mu=1}^{4}p_\mu n_\mu
\end{equation}
also obey the boundary conditions (A.29).

The basic formula, underlying Fourier transformation on the lattice, is (here
$l$ is an integer with $0 \le l \le N - 1$)
\begin{equation}
  \frac{1}{N}\sum_{j=-N/2+1}^{N/2}\exp\!\left(i\,\frac{2\pi}{N}\,lj\right)
  = \frac{1}{N}\sum_{j=0}^{N-1}\exp\!\left(i\,\frac{2\pi}{N}\,lj\right) = \delta_{l0}.
\end{equation}
For $l = 0$ this formula is trivial. For $l \ne 0$ this formula follows from
applying the well-known algebraic identity
\begin{equation}
  \sum_{j=0}^{N-1}q^j = \frac{1 - q^N}{1 - q} \qquad \text{to} \qquad
  q = \exp\!\left(i\,\frac{2\pi}{N}\,l\right).
\end{equation}

We can combine four of the 1D sums in (A.32) to obtain the following
identities:
\begin{equation}
  \frac{1}{|\Lambda|}\sum_{p\in\tilde{\Lambda}}\exp\!\left(i\,p\cdot(n-n')\,a\right)
  = \delta(n-n') = \delta_{n_1n'_1}\delta_{n_2n'_2}\delta_{n_3n'_3}\delta_{n_4n'_4},
\end{equation}
\begin{equation}
  \frac{1}{|\Lambda|}\sum_{n\in\Lambda}\exp\!\left(i\,(p-p')\cdot n\,a\right)
  = \delta(p-p') \equiv \delta_{k_1k'_1}\delta_{k_2k'_2}\delta_{k_3k'_3}\delta_{k_4k'_4}.
\end{equation}
We stress that the right-hand side of (A.35) is a product of four Kronecker
deltas for the integers $k_\mu$ which label the momentum components $p_\mu$
(compare (A.30)).

If we now define the Fourier transform
\begin{equation}
  \tilde{f}(p) = \frac{1}{\sqrt{|\Lambda|}}\sum_{n\in\Lambda}f(n)\,\exp(-i\,p\cdot n\,a),
\end{equation}
we find for the inverse transformation
\begin{equation}
  f(n) = \frac{1}{\sqrt{|\Lambda|}}\sum_{p\in\tilde{\Lambda}}\tilde{f}(p)\,\exp(i\,p\cdot n\,a).
\end{equation}
The last equation follows immediately from inserting the definition of the
Fourier transform in the inverse transformation and using (A.34).

\section{Wilson's formulation of lattice QCD}

In this appendix we collect the defining formulas for Wilson's formulation of
QCD on the lattice. The dynamical variables are the group-valued link
variables $U_\mu(n)$ and the Grassmann-valued fermion fields
$\psi^{(f)}(n)^\alpha_c$, $\bar{\psi}^{(f)}(n)^\alpha_c$. They live on the
links, respectively the sites of our lattice (A.27). Vacuum expectation values
are calculated according to
\begin{equation}
  \ev{O} = \frac{1}{Z}\int D[\psi,\bar{\psi}]\,D[U]\,
  e^{-S_F[\psi,\bar{\psi},U]-S_G[U]}\,O[\psi,\bar{\psi},U],
\end{equation}
where the partition function is given by
\begin{equation}
  Z = \int D[\psi,\bar{\psi}]\,D[U]\,
  e^{-S_F[\psi,\bar{\psi},U]-S_G[U]} .
\end{equation}
The measures over fermion and gauge fields are products over the measures for
the individual field variables:
\begin{align}
  D[\psi,\bar{\psi}] &= \prod_{n\in\Lambda}\prod_{f,\alpha,c}
  d\bar{\psi}^{(f)}(n)^\alpha_c\,d\psi^{(f)}(n)^\alpha_c, \notag\\
  D[U] &= \prod_{n\in\Lambda}\prod_{\mu=1}^{4} dU_\mu(n).
\end{align}
For the individual link variables $U_\mu(n)$ one uses the Haar measure discussed
in Sect.~3.1. For the fermions the rules for Grassmann integration from
Sect.~5.1 apply. The gauge field action for gauge group SU($N$) is given by
\begin{equation}
  S_G[U] = \frac{\beta}{N}\sum_{n\in\Lambda}\sum_{\mu<\nu}\Re\,\tr[1 - U_{\mu\nu}(n)],
\end{equation}
where the plaquettes are defined as
\begin{equation}
  \begin{split}
    U_{\mu\nu}(n) &= U_\mu(n)\,U_\nu(n+\hat{\mu})\,U_{-\mu}(n+\hat{\mu}+\hat{\nu})\,U_{-\nu}(n+\hat{\nu})\\
    &= U_\mu(n)\,U_\nu(n+\hat{\mu})\,U_\mu(n+\hat{\nu})^\dagger\,U_\nu(n)^\dagger .
  \end{split}
\end{equation}
The fermion action is a sum over $N_f$ flavors:
\begin{equation}
  S_F[\psi,\bar{\psi},U] = \sum_{f=1}^{N_f}a^4\sum_{n,m\in\Lambda}
  \bar{\psi}^{(f)}(n)\,D^{(f)}(n|m)\,\psi^{(f)}(m)
\end{equation}
and the lattice Dirac operator is given by
\begin{equation}
  D^{(f)}(n|m)^{\alpha\beta}_{ab} = \left( m^{(f)} + \frac{4}{a} \right)
  \delta_{\alpha\beta}\delta_{ab}\delta_{n,m}
  - \frac{1}{2a}\sum_{\mu=\pm1}^{\pm4}(1-\gamma_\mu)_{\alpha\beta}\,
  U_\mu(n)_{ab}\,\delta_{n+\hat{\mu},m}.
\end{equation}
In the plaquette definition and in the last equation we use the conventions
\begin{equation}
  \gamma_{-\mu} = -\gamma_\mu,\qquad
  U_{-\mu}(n) = U_\mu(n-\hat{\mu})^\dagger,\qquad \mu = 1, 2, 3, 4 .
\end{equation}
We remark that Wilson's Dirac operator is $\gamma_5$-hermitian, i.e., it obeys
\begin{equation}
  \gamma_5 D \gamma_5 = D^\dagger .
\end{equation}

\section{A few formulas for matrix algebra}

In quantum mechanics one usually deals with hermitian or unitary matrices,
while in lattice QCD often more general matrices occur. In this appendix we
list a few results for general complex matrices together with short remarks
concerning their proof (for a more detailed account see, e.g., [5]).

The basic result for general complex matrices is that they are unitarily
equivalent to upper triangular matrices: Let $M$ be a complex-valued
$N\times N$ matrix. Then there exists a unitary matrix $U$ and an upper
triangular matrix $T$, such that
\begin{equation}
  U^\dagger M U = T .
\end{equation}
This result can be proven by induction in $N$. The elements $t_j$ on the
diagonal of $T$ are the roots of the characteristic polynomial of $M$ since
\begin{equation}
  P(\lambda) = \det[M - \lambda\,\mathbb{1}] = \det[T - \lambda\,\mathbb{1}]
  = \prod_{j=1}^{N}(t_j - \lambda).
\end{equation}

An important consequence of this result is a unique classification of matrices
that can be diagonalized with a unitary transformation. A complex matrix $M$
is called normal if it commutes with its hermitian conjugate, i.e.,
$[M, M^\dagger] = 0$. It is obvious that hermitian or unitary matrices are
normal. The announced result is: If and only if $M$ is normal, then there
exists a unitary matrix $U$ such that
\begin{equation}
  U^\dagger M U = D ,
\end{equation}
where $D$ is diagonal. It is straightforward to see that a matrix $M$ which is
unitarily equivalent to a diagonal matrix is normal. To prove the other
direction we first note that the normality of $M$ implies the normality of the
triangular matrix $T$ corresponding to $M$. When evaluating explicitly the two
sides of the normality condition, $T^\dagger T = T T^\dagger$, for the upper
triangular matrix $T$, one concludes that $T$ must be diagonal and the
statement is proven. Equations (A.49) and (A.48) imply that a normal matrix
has a complete set of orthonormal eigenvectors, the columns of $U$.

The existence of a complete orthonormal set of eigenvectors $v^{(j)}$ with
eigenvalues $\lambda^{(j)}$ can be used to represent the matrix $M$ in the
form
\begin{equation}
  M = \sum_{j=1}^{N}\lambda^{(j)}\,v^{(j)}v^{(j)\dagger},
\end{equation}
the so-called spectral representation. On the right-hand side of this equation
matrix/vector notation was used to write the dyadic product
$v^{(j)}v^{(j)\dagger}$. The spectral representation of the matrix can be used
to define a function of $M$ in terms of the function for the eigenvalues, if
this exists. This gives rise to the spectral theorem
\begin{equation}
  f(M) = \sum_{j=1}^{N}f\!\left(\lambda^{(j)}\right) v^{(j)}v^{(j)\dagger}.
\end{equation}

We finally discuss a formula for the expansion of the determinant:
\begin{equation}
  \det[\mathbb{1} - M] = \exp\!\left(\tr[\ln(\mathbb{1} - M)]\right).
\end{equation}
In this equation $M$ is a complex matrix and the logarithm (where it exists)
is defined through its series expansion. The proof of this result applies the
triangularization (A.47):
\begin{align}
  \det[\mathbb{1} - M] &= \det[\mathbb{1} - T] = \prod_{j=1}^{N}(1 - t_j)
  = \exp\!\left(\sum_{j=1}^{N}\ln(1 - t_j)\right) \notag\\
  &= \exp\!\left(-\sum_{j=1}^{N}\sum_{n=1}^{\infty}\frac{1}{n}(t_j)^n\right)
  = \exp\!\left(-\sum_{n=1}^{\infty}\frac{1}{n}\tr[T^n]\right) \notag\\
  &= \exp\!\left(\tr[\ln(\mathbb{1} - M)]\right).
\end{align}
In the fifth step we have used the fact that when evaluating powers of a
triangular matrix the diagonal elements do not mix with other entries of the
matrix. In the last step we used $\tr[T^n] = \tr[M^n]$ which follows from
(A.47). Since a matrix $A$ may always be written as $A = \mathbb{1} - M$, the
result is often stated as
\begin{equation}
  \det[A] = \exp\!\left(\tr[\ln A]\right).
\end{equation}

\section*{References}

\begin{enumerate}[label={[\arabic*]}]
  \item H. Georgi: Lie Algebras in Particle Physics (Benjamin/Cummings,
    Reading, Massachusetts 1982)
  \item H. F. Jones: Groups, Representations and Physics (Hilger, Bristol 1990)
  \item M. Hamermesh: Group Theory and Its Application to Physical Problems
    (Addison-Wesley, Reading, Massachusetts 1964)
  \item M. E. Peskin and D. V. Schroeder: An Introduction to Quantum Field
    Theory (Addison-Wesley, Reading, Massachusetts 1995)
  \item P. Lancaster: Theory of Matrices (Academic Press, New York 1969)
\end{enumerate}
, so the chapter is
%  automatically lettered "A" and the sections become A.1, A.2, ...
% =====================================================================

\chapter{Appendix}

\section{The Lie groups SU($N$)}

In this appendix we collect basic definitions and conventions for the Lie
groups SU($N$) -- the special unitary groups -- and the corresponding Lie
algebras $su(N)$. For a more detailed presentation we refer the reader to
[1--3].

\subsection{Basic properties}

The defining representation of SU($N$) is given by complex $N\times N$
matrices which are unitary and have determinant 1. This set of matrices is
closed under matrix multiplication: Let $\Omega_1$ and $\Omega_2$ be elements
of SU($N$), i.e., they obey $\Omega_i^\dagger = \Omega_i^{-1}$ and
$\det[\Omega_i] = 1$. Using standard linear algebra manipulations we obtain
\begin{equation}
  \begin{gathered}
    (\Omega_1\Omega_2)^\dagger = \Omega_2^\dagger\Omega_1^\dagger
    = \Omega_2^{-1}\Omega_1^{-1} = (\Omega_1\Omega_2)^{-1},\\[2pt]
    \det[\Omega_1\Omega_2] = \det[\Omega_1]\det[\Omega_2] = 1
  \end{gathered}
\end{equation}
and have thus established that also the product of two SU($N$) matrices is an
SU($N$) matrix. The unit matrix is also in SU($N$) and for each matrix in
SU($N$) there exists an inverse (the hermitian conjugate matrix). Thus, the
set SU($N$) forms a group. Since the group operation -- the matrix
multiplication -- is not commutative the groups SU($N$) are non-abelian
groups.

\subsection{Lie algebra}

Let us now count how many real parameters are needed to describe the matrices
in SU($N$). A complex $N\times N$ matrix has $2N^2$ real parameters. The
requirement of unitarity introduces $N^2$ independent conditions which the
parameters have to obey. One more parameter is used for obeying the determinant
condition such that one needs a total of $N^2-1$ real parameters for describing
SU($N$) matrices.

A convenient way of representing SU($N$) matrices is to write them as
exponentials of basis matrices $T_j$, the so-called generators. In particular,
we write an element $\Omega$ of SU($N$) as
\begin{equation}
  \Omega = \exp\!\left(i\sum_{j=1}^{N^2-1}\omega^{(j)} T_j\right),
\end{equation}
where $\omega^{(j)}$, $j = 1, 2, \ldots, N^2-1$, are the real numbers needed to
parameterize $\Omega$. We remark that the parameters $\omega^{(j)}$ can be
changed continuously, making SU($N$) so-called Lie groups that are groups whose
elements depend continuously on their parameters. In order to cover all of the
group space, the parameters have to be varied only over finite intervals,
making the groups SU($N$) so-called compact Lie groups.

The generators $T_j$, $j = 1, 2, \ldots, N^2-1$, are chosen as traceless,
complex, and hermitian $N\times N$ matrices obeying the normalization condition
\begin{equation}
  \tr[T_j T_k] = \half\,\delta_{jk}.
\end{equation}
In addition, they are related among each other by an algebra of commutation
relations
\begin{equation}
  [T_j, T_k] = i\, f_{jkl}\, T_l .
\end{equation}
The completely anti-symmetric coefficients $f_{jkl}$ are the so-called
structure constants. Below we will give an explicit representation of the
generators for the groups SU(2) and SU(3).

Let us verify that the representation of $\Omega$ indeed describes elements of
SU($N$). Using the facts that the generators are hermitian and that the
$\omega^{(j)}$ are real, one finds that hermitian conjugation of the right-hand
side of the exponent simply produces an extra minus sign in the exponent (from
the complex conjugation of $i$). Thus, it is obvious that $\Omega^\dagger =
\Omega^{-1}$. To show that the determinant equals 1, we use the equation
\begin{equation}
  \det\Omega = \exp(\tr[\ln\Omega])
  = \exp\!\left(i\sum_{j=1}^{N^2-1}\omega^{(j)}\tr T_j\right) = e^0 = 1,
\end{equation}
where in first step we have used a formula for the determinant (see (A.54)
below) and in the third step we have used the fact that the $T_j$ are
traceless.

Not only the group elements but also the exponents of our representation have
an interesting structure. The linear combinations
\begin{equation}
  \sum_{j=1}^{N^2-1}\omega^{(j)} T_j
\end{equation}
of the $T_j$ form the so-called Lie algebra $su(N)$. Their commutation
properties are governed by the commutation relations. Elements of $su(N)$ are
also complex $N\times N$ matrices but have properties different from the
elements of the group. One important difference is the fact that the unit
matrix is contained in the group (all $\omega^{(j)} = 0$), while it is not an
element of the algebra (all $T_j$ are traceless).

\subsection{Generators for SU(2) and SU(3)}

The standard representation of the generators for SU(2) is given by
\begin{equation}
  T_j = \half\,\sigma_j,
\end{equation}
with the Pauli matrices
\begin{equation}
  \sigma_1 = \begin{pmatrix} 0 & 1 \\ 1 & 0 \end{pmatrix},\qquad
  \sigma_2 = \begin{pmatrix} 0 & -i \\ i & 0 \end{pmatrix},\qquad
  \sigma_3 = \begin{pmatrix} 1 & 0 \\ 0 & -1 \end{pmatrix}.
\end{equation}
In this case the structure constants are particularly simple, given by the
completely anti-symmetric tensor, i.e., $f_{jkl} = \varepsilon_{jkl}$.

For SU(3) the generators are given by
\begin{equation}
  T_j = \half\,\lambda_j .
\end{equation}
The Gell--Mann matrices $\lambda_j$ are $3\times 3$ generalizations of the
Pauli matrices:
\begin{equation}
  \begin{gathered}
    \lambda_1 = \begin{pmatrix} 0 & 1 & 0 \\ 1 & 0 & 0 \\ 0 & 0 & 0 \end{pmatrix},\qquad
    \lambda_2 = \begin{pmatrix} 0 & -i & 0 \\ i & 0 & 0 \\ 0 & 0 & 0 \end{pmatrix},\qquad
    \lambda_3 = \begin{pmatrix} 1 & 0 & 0 \\ 0 & -1 & 0 \\ 0 & 0 & 0 \end{pmatrix},\\[4pt]
    \lambda_4 = \begin{pmatrix} 0 & 0 & 1 \\ 0 & 0 & 0 \\ 1 & 0 & 0 \end{pmatrix},\qquad
    \lambda_5 = \begin{pmatrix} 0 & 0 & -i \\ 0 & 0 & 0 \\ i & 0 & 0 \end{pmatrix},\qquad
    \lambda_6 = \begin{pmatrix} 0 & 0 & 0 \\ 0 & 0 & 1 \\ 0 & 1 & 0 \end{pmatrix},\\[4pt]
    \lambda_7 = \begin{pmatrix} 0 & 0 & 0 \\ 0 & 0 & -i \\ 0 & i & 0 \end{pmatrix},\qquad
    \lambda_8 = \frac{1}{\sqrt{3}}
    \begin{pmatrix} 1 & 0 & 0 \\ 0 & 1 & 0 \\ 0 & 0 & -2 \end{pmatrix}.
  \end{gathered}
\end{equation}

\subsection{Derivatives of group elements}

Let us now show an important property of derivatives of group elements. If
$\Omega(\omega)$ is an element of SU($N$) then
\begin{equation}
  M_k(\omega) = i\,\frac{\partial\Omega(\omega)}{\partial\omega^{(k)}}\,
  \Omega(\omega)^\dagger \in su(N),
\end{equation}
i.e., the derivative times the hermitian conjugate is in the Lie algebra. In
order to prove this statement we have to show the defining properties of Lie
algebra elements, i.e., we must show that $M_k(\omega)$ is hermitian and
traceless.

Showing the hermiticity of $M_k(\omega)$ is straightforward. By differentiating
$\Omega(\omega)\Omega(\omega)^\dagger = 1$ with respect to $\omega^{(k)}$ one
finds
\begin{equation}
  \frac{\partial\Omega(\omega)}{\partial\omega^{(k)}}\,\Omega(\omega)^\dagger
  + \Omega(\omega)\,\frac{\partial\Omega(\omega)^\dagger}{\partial\omega^{(k)}} = 0.
\end{equation}
Thus
\begin{align}
  M_k(\omega)^\dagger &=
  \left[ i\,\frac{\partial\Omega(\omega)}{\partial\omega^{(k)}}\,
  \Omega(\omega)^\dagger \right]^\dagger
  = -i\,\Omega(\omega)\,\frac{\partial\Omega(\omega)^\dagger}{\partial\omega^{(k)}} \notag\\
  &= i\,\frac{\partial\Omega(\omega)}{\partial\omega^{(k)}}\,
  \Omega(\omega)^\dagger = M_k(\omega),
\end{align}
where we used (A.12) in the third step.

In order to show that $M_k(\omega)$ is traceless we use the fact that the
determinant of an SU($N$) matrix equals 1 and we differentiate
$\det[\Omega(\omega)]$ with respect to $\omega^{(k)}$. We obtain
\begin{align}
  0 &= \frac{\partial\det[\Omega(\omega)]}{\partial\omega^{(k)}}
  = \frac{\partial\det[\Omega(\omega)]}{\partial\Omega(\omega)_{ab}}\,
  \frac{\partial\Omega(\omega)_{ab}}{\partial\omega^{(k)}} \notag\\
  &= \det[\Omega(\omega)]\left(\Omega(\omega)^{-1}\right)_{ba}\,
  \frac{\partial\Omega(\omega)_{ab}}{\partial\omega^{(k)}}
  = \tr\!\left[\frac{\partial\Omega(\omega)}{\partial\omega^{(k)}}\,
  \Omega(\omega)^\dagger\right],
\end{align}
where in the first step we applied the chain rule for derivatives. In the second
step we used a standard formula for the derivative of the determinant
$\det[\Omega]$ with respect to an entry $\Omega_{ab}$ of the matrix $\Omega$.
In the third step we used $\det[\Omega] = 1$ and $\Omega^{-1} = \Omega^\dagger$.
Equation (A.14) establishes that $M_k(\omega)$ is also traceless and thus we
have shown $M_k(\omega)\in su(N)$.

From (A.11) it follows that for the gauge transformation matrices $\Omega(x)$
with coefficients $\omega^{(k)}(x)$, depending on the space--time coordinate
$x$, the combination $i(\partial_\mu\Omega(x))\Omega(x)^\dagger$ is in the Lie
algebra, since
\begin{equation}
  i(\partial_\mu\Omega(x))\,\Omega(x)^\dagger
  = \sum_k i\,\frac{\partial\Omega(\omega(x))}{\partial\omega^{(k)}(x)}\,
  \Omega(\omega(x))^\dagger\,\partial_\mu\omega^{(k)}(x),
\end{equation}
and the right-hand side is a linear combination of $su(N)$ elements with real
coefficients $\partial_\mu\omega^{(k)}(x)$.

\section{Gamma matrices}

The Euclidean gamma matrices $\gamma_\mu$, $\mu = 1, 2, 3, 4$ can be
constructed from the Minkowski gamma matrices $\gamma_\mu^M$, $\mu = 0, 1, 2,
3$. The latter obey
\begin{equation}
  \{\gamma_\mu^M, \gamma_\nu^M\} = 2\,g_{\mu\nu}\,\mathbb{1},
\end{equation}
with the metric tensor given by $g_{\mu\nu} = \mathrm{diag}(1,-1,-1,-1)$ and
$\mathbb{1}$ is the $4\times 4$ unit matrix. Thus when we define the Euclidean
matrices $\gamma_\mu$ by setting
\begin{equation}
  \gamma_1 = -i\gamma_1^M,\quad
  \gamma_2 = -i\gamma_2^M,\quad
  \gamma_3 = -i\gamma_3^M,\quad
  \gamma_4 = \gamma_0^M,
\end{equation}
we obtain the Euclidean anti-commutation relations
\begin{equation}
  \{\gamma_\mu, \gamma_\nu\} = 2\,\delta_{\mu\nu}\,\mathbb{1}.
\end{equation}

In addition to the matrices $\gamma_\mu$, $\mu = 1, 2, 3, 4$ we define the
matrix $\gamma_5$ as the product
\begin{equation}
  \gamma_5 = \gamma_1\gamma_2\gamma_3\gamma_4 .
\end{equation}
The matrix $\gamma_5$ anti-commutes with all other gamma matrices $\gamma_\mu$,
$\mu = 1, 2, 3, 4$ and obeys $\gamma_5^2 = 1$.

An explicit representation of the Euclidean gamma matrices can be obtained from
a representation of the Minkowski gamma matrices (see, e.g., [4]). Here we give
the so-called chiral representation where $\gamma_5$ (the chirality operator)
is diagonal:
\begin{equation}
  \gamma_{1,2,3} = \begin{pmatrix} 0 & -i\sigma_{1,2,3} \\ i\sigma_{1,2,3} & 0 \end{pmatrix},\qquad
  \gamma_4 = \begin{pmatrix} 0 & \mathbb{1}_2 \\ \mathbb{1}_2 & 0 \end{pmatrix},\qquad
  \gamma_5 = \begin{pmatrix} \mathbb{1}_2 & 0 \\ 0 & -\mathbb{1}_2 \end{pmatrix},
\end{equation}
where the $\sigma_j$ are the Pauli matrices and $\mathbb{1}_2$ is the
$2\times 2$ unit matrix. More explicitly the Euclidean gamma matrices read
\begin{equation}
  \begin{gathered}
    \gamma_1 = \begin{pmatrix} 0 & 0 & 0 & -i \\ 0 & 0 & -i & 0 \\ 0 & i & 0 & 0 \\ i & 0 & 0 & 0 \end{pmatrix},\qquad
    \gamma_2 = \begin{pmatrix} 0 & 0 & 0 & -1 \\ 0 & 0 & 1 & 0 \\ 0 & 1 & 0 & 0 \\ -1 & 0 & 0 & 0 \end{pmatrix},\\[4pt]
    \gamma_3 = \begin{pmatrix} 0 & 0 & -i & 0 \\ 0 & 0 & 0 & i \\ i & 0 & 0 & 0 \\ 0 & -i & 0 & 0 \end{pmatrix},\qquad
    \gamma_4 = \begin{pmatrix} 0 & 0 & 1 & 0 \\ 0 & 0 & 0 & 1 \\ 1 & 0 & 0 & 0 \\ 0 & 1 & 0 & 0 \end{pmatrix},\\[4pt]
    \gamma_5 = \begin{pmatrix} 1 & 0 & 0 & 0 \\ 0 & 1 & 0 & 0 \\ 0 & 0 & -1 & 0 \\ 0 & 0 & 0 & -1 \end{pmatrix}.
  \end{gathered}
\end{equation}
In addition to the anti-commutation relation (A.18) the gamma matrices obey
(here $\mu = 1, \ldots, 5$)
\begin{equation}
  \gamma_\mu = \gamma_\mu^\dagger = \gamma_\mu^{-1}.
\end{equation}

When we discuss charge conjugation, we need the charge conjugation matrix $C$
defined through the relations ($\mu = 1, \ldots, 4$)
\begin{equation}
  C\gamma_\mu C^{-1} = -\gamma_\mu^T .
\end{equation}
Using the explicit form of the gamma matrices it is easy to see that in the
chiral representation the charge conjugation matrix is given by
\begin{equation}
  C = i\gamma_2\gamma_4 .
\end{equation}
It obeys
\begin{equation}
  C = C^{-1} = C^\dagger = -C^T .
\end{equation}

We finally quote a simple formula for the inverse of linear combinations of
gamma matrices ($a, b_\mu \in \mathbb{R}$):
\begin{equation}
  \left( a\,\mathbb{1} + i\sum_{\mu=1}^{4}\gamma_\mu b_\mu \right)^{-1}
  = \frac{ a\,\mathbb{1} - i\sum_{\mu=1}^{4}\gamma_\mu b_\mu }{ a^2 + \sum_{\mu=1}^{4}b_\mu^2 }.
\end{equation}
This formula can be verified by multiplying both sides with
$a\mathbb{1} + i\sum_\mu \gamma_\mu b_\mu$.

\section{Fourier transformation on the lattice}

The goal of this appendix is to discuss the Fourier transform $\tilde{f}(p)$
of functions $f(n)$ defined on the lattice $\Lambda$. The lattice is given by
\begin{equation}
  \Lambda = \left\{ n = (n_1, n_2, n_3, n_4) \,\big|\, n_\mu = 0, 1, \ldots, N_\mu - 1 \right\},
\end{equation}
and in most of our applications we have $N_1 = N_2 = N_3 = N$, $N_4 = N_T$.
For the total number of lattice points we introduce the abbreviation
\begin{equation}
  |\Lambda| = N_1 N_2 N_3 N_4 .
\end{equation}

We impose toroidal boundary conditions
\begin{equation}
  f(n + \hat{\mu}N_\mu) = e^{i 2\pi\theta_\mu}\, f(n)
\end{equation}
for each of the directions $\mu$. Here $\hat{\mu}$ denotes the unit vector in
$\mu$-direction. Directions with periodic boundary conditions have
$\theta_\mu = 0$, anti-periodic boundary conditions correspond to
$\theta_\mu = 1/2$.

The momentum space $\tilde{\Lambda}$, which corresponds to the lattice
$\Lambda$ with the boundary conditions (A.29), is defined as
\begin{equation}
  \tilde{\Lambda} = \left\{ p = (p_1, p_2, p_3, p_4) \,\Big|\, p_\mu =
  \frac{2\pi}{aN_\mu}(k_\mu + \theta_\mu),\; k_\mu = -\frac{N_\mu}{2} + 1, \ldots, \frac{N_\mu}{2} \right\}.
\end{equation}
The boundary phases $\theta_\mu$ have to be included in the definition of the
momenta $p_\mu$ such that the plane waves
\begin{equation}
  \exp(i\,p\cdot n\,a) \qquad \text{with} \qquad p\cdot n = \sum_{\mu=1}^{4}p_\mu n_\mu
\end{equation}
also obey the boundary conditions (A.29).

The basic formula, underlying Fourier transformation on the lattice, is (here
$l$ is an integer with $0 \le l \le N - 1$)
\begin{equation}
  \frac{1}{N}\sum_{j=-N/2+1}^{N/2}\exp\!\left(i\,\frac{2\pi}{N}\,lj\right)
  = \frac{1}{N}\sum_{j=0}^{N-1}\exp\!\left(i\,\frac{2\pi}{N}\,lj\right) = \delta_{l0}.
\end{equation}
For $l = 0$ this formula is trivial. For $l \ne 0$ this formula follows from
applying the well-known algebraic identity
\begin{equation}
  \sum_{j=0}^{N-1}q^j = \frac{1 - q^N}{1 - q} \qquad \text{to} \qquad
  q = \exp\!\left(i\,\frac{2\pi}{N}\,l\right).
\end{equation}

We can combine four of the 1D sums in (A.32) to obtain the following
identities:
\begin{equation}
  \frac{1}{|\Lambda|}\sum_{p\in\tilde{\Lambda}}\exp\!\left(i\,p\cdot(n-n')\,a\right)
  = \delta(n-n') = \delta_{n_1n'_1}\delta_{n_2n'_2}\delta_{n_3n'_3}\delta_{n_4n'_4},
\end{equation}
\begin{equation}
  \frac{1}{|\Lambda|}\sum_{n\in\Lambda}\exp\!\left(i\,(p-p')\cdot n\,a\right)
  = \delta(p-p') \equiv \delta_{k_1k'_1}\delta_{k_2k'_2}\delta_{k_3k'_3}\delta_{k_4k'_4}.
\end{equation}
We stress that the right-hand side of (A.35) is a product of four Kronecker
deltas for the integers $k_\mu$ which label the momentum components $p_\mu$
(compare (A.30)).

If we now define the Fourier transform
\begin{equation}
  \tilde{f}(p) = \frac{1}{\sqrt{|\Lambda|}}\sum_{n\in\Lambda}f(n)\,\exp(-i\,p\cdot n\,a),
\end{equation}
we find for the inverse transformation
\begin{equation}
  f(n) = \frac{1}{\sqrt{|\Lambda|}}\sum_{p\in\tilde{\Lambda}}\tilde{f}(p)\,\exp(i\,p\cdot n\,a).
\end{equation}
The last equation follows immediately from inserting the definition of the
Fourier transform in the inverse transformation and using (A.34).

\section{Wilson's formulation of lattice QCD}

In this appendix we collect the defining formulas for Wilson's formulation of
QCD on the lattice. The dynamical variables are the group-valued link
variables $U_\mu(n)$ and the Grassmann-valued fermion fields
$\psi^{(f)}(n)^\alpha_c$, $\bar{\psi}^{(f)}(n)^\alpha_c$. They live on the
links, respectively the sites of our lattice (A.27). Vacuum expectation values
are calculated according to
\begin{equation}
  \ev{O} = \frac{1}{Z}\int D[\psi,\bar{\psi}]\,D[U]\,
  e^{-S_F[\psi,\bar{\psi},U]-S_G[U]}\,O[\psi,\bar{\psi},U],
\end{equation}
where the partition function is given by
\begin{equation}
  Z = \int D[\psi,\bar{\psi}]\,D[U]\,
  e^{-S_F[\psi,\bar{\psi},U]-S_G[U]} .
\end{equation}
The measures over fermion and gauge fields are products over the measures for
the individual field variables:
\begin{align}
  D[\psi,\bar{\psi}] &= \prod_{n\in\Lambda}\prod_{f,\alpha,c}
  d\bar{\psi}^{(f)}(n)^\alpha_c\,d\psi^{(f)}(n)^\alpha_c, \notag\\
  D[U] &= \prod_{n\in\Lambda}\prod_{\mu=1}^{4} dU_\mu(n).
\end{align}
For the individual link variables $U_\mu(n)$ one uses the Haar measure discussed
in Sect.~3.1. For the fermions the rules for Grassmann integration from
Sect.~5.1 apply. The gauge field action for gauge group SU($N$) is given by
\begin{equation}
  S_G[U] = \frac{\beta}{N}\sum_{n\in\Lambda}\sum_{\mu<\nu}\Re\,\tr[1 - U_{\mu\nu}(n)],
\end{equation}
where the plaquettes are defined as
\begin{equation}
  \begin{split}
    U_{\mu\nu}(n) &= U_\mu(n)\,U_\nu(n+\hat{\mu})\,U_{-\mu}(n+\hat{\mu}+\hat{\nu})\,U_{-\nu}(n+\hat{\nu})\\
    &= U_\mu(n)\,U_\nu(n+\hat{\mu})\,U_\mu(n+\hat{\nu})^\dagger\,U_\nu(n)^\dagger .
  \end{split}
\end{equation}
The fermion action is a sum over $N_f$ flavors:
\begin{equation}
  S_F[\psi,\bar{\psi},U] = \sum_{f=1}^{N_f}a^4\sum_{n,m\in\Lambda}
  \bar{\psi}^{(f)}(n)\,D^{(f)}(n|m)\,\psi^{(f)}(m)
\end{equation}
and the lattice Dirac operator is given by
\begin{equation}
  D^{(f)}(n|m)^{\alpha\beta}_{ab} = \left( m^{(f)} + \frac{4}{a} \right)
  \delta_{\alpha\beta}\delta_{ab}\delta_{n,m}
  - \frac{1}{2a}\sum_{\mu=\pm1}^{\pm4}(1-\gamma_\mu)_{\alpha\beta}\,
  U_\mu(n)_{ab}\,\delta_{n+\hat{\mu},m}.
\end{equation}
In the plaquette definition and in the last equation we use the conventions
\begin{equation}
  \gamma_{-\mu} = -\gamma_\mu,\qquad
  U_{-\mu}(n) = U_\mu(n-\hat{\mu})^\dagger,\qquad \mu = 1, 2, 3, 4 .
\end{equation}
We remark that Wilson's Dirac operator is $\gamma_5$-hermitian, i.e., it obeys
\begin{equation}
  \gamma_5 D \gamma_5 = D^\dagger .
\end{equation}

\section{A few formulas for matrix algebra}

In quantum mechanics one usually deals with hermitian or unitary matrices,
while in lattice QCD often more general matrices occur. In this appendix we
list a few results for general complex matrices together with short remarks
concerning their proof (for a more detailed account see, e.g., [5]).

The basic result for general complex matrices is that they are unitarily
equivalent to upper triangular matrices: Let $M$ be a complex-valued
$N\times N$ matrix. Then there exists a unitary matrix $U$ and an upper
triangular matrix $T$, such that
\begin{equation}
  U^\dagger M U = T .
\end{equation}
This result can be proven by induction in $N$. The elements $t_j$ on the
diagonal of $T$ are the roots of the characteristic polynomial of $M$ since
\begin{equation}
  P(\lambda) = \det[M - \lambda\,\mathbb{1}] = \det[T - \lambda\,\mathbb{1}]
  = \prod_{j=1}^{N}(t_j - \lambda).
\end{equation}

An important consequence of this result is a unique classification of matrices
that can be diagonalized with a unitary transformation. A complex matrix $M$
is called normal if it commutes with its hermitian conjugate, i.e.,
$[M, M^\dagger] = 0$. It is obvious that hermitian or unitary matrices are
normal. The announced result is: If and only if $M$ is normal, then there
exists a unitary matrix $U$ such that
\begin{equation}
  U^\dagger M U = D ,
\end{equation}
where $D$ is diagonal. It is straightforward to see that a matrix $M$ which is
unitarily equivalent to a diagonal matrix is normal. To prove the other
direction we first note that the normality of $M$ implies the normality of the
triangular matrix $T$ corresponding to $M$. When evaluating explicitly the two
sides of the normality condition, $T^\dagger T = T T^\dagger$, for the upper
triangular matrix $T$, one concludes that $T$ must be diagonal and the
statement is proven. Equations (A.49) and (A.48) imply that a normal matrix
has a complete set of orthonormal eigenvectors, the columns of $U$.

The existence of a complete orthonormal set of eigenvectors $v^{(j)}$ with
eigenvalues $\lambda^{(j)}$ can be used to represent the matrix $M$ in the
form
\begin{equation}
  M = \sum_{j=1}^{N}\lambda^{(j)}\,v^{(j)}v^{(j)\dagger},
\end{equation}
the so-called spectral representation. On the right-hand side of this equation
matrix/vector notation was used to write the dyadic product
$v^{(j)}v^{(j)\dagger}$. The spectral representation of the matrix can be used
to define a function of $M$ in terms of the function for the eigenvalues, if
this exists. This gives rise to the spectral theorem
\begin{equation}
  f(M) = \sum_{j=1}^{N}f\!\left(\lambda^{(j)}\right) v^{(j)}v^{(j)\dagger}.
\end{equation}

We finally discuss a formula for the expansion of the determinant:
\begin{equation}
  \det[\mathbb{1} - M] = \exp\!\left(\tr[\ln(\mathbb{1} - M)]\right).
\end{equation}
In this equation $M$ is a complex matrix and the logarithm (where it exists)
is defined through its series expansion. The proof of this result applies the
triangularization (A.47):
\begin{align}
  \det[\mathbb{1} - M] &= \det[\mathbb{1} - T] = \prod_{j=1}^{N}(1 - t_j)
  = \exp\!\left(\sum_{j=1}^{N}\ln(1 - t_j)\right) \notag\\
  &= \exp\!\left(-\sum_{j=1}^{N}\sum_{n=1}^{\infty}\frac{1}{n}(t_j)^n\right)
  = \exp\!\left(-\sum_{n=1}^{\infty}\frac{1}{n}\tr[T^n]\right) \notag\\
  &= \exp\!\left(\tr[\ln(\mathbb{1} - M)]\right).
\end{align}
In the fifth step we have used the fact that when evaluating powers of a
triangular matrix the diagonal elements do not mix with other entries of the
matrix. In the last step we used $\tr[T^n] = \tr[M^n]$ which follows from
(A.47). Since a matrix $A$ may always be written as $A = \mathbb{1} - M$, the
result is often stated as
\begin{equation}
  \det[A] = \exp\!\left(\tr[\ln A]\right).
\end{equation}

\section*{References}

\begin{enumerate}[label={[\arabic*]}]
  \item H. Georgi: Lie Algebras in Particle Physics (Benjamin/Cummings,
    Reading, Massachusetts 1982)
  \item H. F. Jones: Groups, Representations and Physics (Hilger, Bristol 1990)
  \item M. Hamermesh: Group Theory and Its Application to Physical Problems
    (Addison-Wesley, Reading, Massachusetts 1964)
  \item M. E. Peskin and D. V. Schroeder: An Introduction to Quantum Field
    Theory (Addison-Wesley, Reading, Massachusetts 1995)
  \item P. Lancaster: Theory of Matrices (Academic Press, New York 1969)
\end{enumerate}
